% vim:fileencoding=utf-8
% -*- encoding: utf-8 -*-
\documentclass[11pt,norsk,a4paper,pointlessnumbers]{scrbook}
\usepackage[utf8]{inputenc}
\usepackage[T1]{fontenc}
\usepackage{babel}
\usepackage{graphicx}
\usepackage{varioref}
\usepackage{amsmath,url}
%\usepackage{epsfig}
\usepackage{moreverb}
\usepackage{float}
\usepackage{fancyheadings}

\title{Hvitebok\\
       \quad \\
       \quad \\
       \includegraphics[height=8cm]{rflogo} }
\author{Realistforeningen}

\newcommand{\clentry}[4]{\section*{Revisjonsnummer #1}\begin{tabular}{ll}\textbf{Tid} & #3 \\ \textbf{Lagt inn av} & #2 \end{tabular}}

\renewcommand{\labelitemi}{--}
\renewcommand{\labelitemii}{$\bullet$}

\begin{document}

\frontmatter

\maketitle{}
\tableofcontents{}


\chapter{Forord}
Denne hviteboken er en samling med tekster av forskjellig slag som har
som felles mål å overføre erfaringer og tradisjoner i Realistforeningen
videre til neste generasjon av aktive studenter. Den er på ingen måte
ment å være en utfyllende og allvitende bok. Nytten av boken står og
faller på at innholdet holdes aktuelt, og det er derfor viktig at den
oppdateres kontinuerlig med de endringer og erfaringer som kommer til.

Boken kan brukes på flere måter, både som en innføring for utenforstående
i hva Realistforeningen holder på med, som opplæringsmateriale for nye
personer i aktive verv, og til å klarere definere ansvarsområder.
Rekkefølgen på de forskjellige delene er mer eller mindre tilfeldig,
og de fleste vil nok bruke den mest som et oppslagsverk.

Jeg håper dette arbeidet vi har lagt ned gir en god hjelp til alle
som kommer etter oss, og at boka vil bli videreutviklet til å bli
langt mer etter hvert.

Martin Alnæs, Formann høsten 2003


\mainmatter

\chapter{Generelt}
Her følger noen enkle betraktninger om det å drive en studentforeninger,
som forhåpentligvis kan gjøre jobben lettere for nye personer.
\section{Kommunikasjon}
Den vanligste �rsaken til at noe ikke g�r som planlagt, er
kommunikasjonsproblemer. Noen enkle tips for � unng� dette:

En kort samtale ansikt til ansikt eller over telefon er ofte raskere og klarere enn tekstmeldinger eller e-post.

\subsection{E-post}
Dette kan se ut som motsetninger, men det reflekterer at bruk av e-post
er et tveegget sverd. Brukt p� riktig m�te er det veldig effektivt,
brukt p� feil m�te skaper det kaos.
\begin{itemize}
\item Bruk e-post aktivt, b�de personlig og lister. Les e-posten din daglig. Dette kan lette arbeidet utrolig mye i mange tilfeller.
\item Ikke skriv for mye til listene, det undergraver effektiviteten fordi mange da slutter � lese s� n�ye det som kommer.
\item Aldri anta at en e-post har blir lest. Hvis den er viktig har den sannsynligvis ikke det.
\item E-post fungerer generelt ikke til hastemeldinger.
\item Bruk riktig liste til meldingene dine. Svada h�rer ikke hjemme p� seri�se lister.
\item Visste du at du kan sortere e-post fra lister automatisk i forskjellige mapper? Anbefales! Sp�r noen om hjelp hvis du ikke vet hvordan.
\end{itemize}


\section{Ansvar}
Jo lenger opp i rekkene du kommer, jo oftere m� du avst� fra � gj�re enkelt oppgaver for � ikke slite deg ut. Husk at ansvar for at noe blir gjennomf�rt ikke betyr at du m� gj�re det selv. Husk ogs� at ansvar for noe ogs� inneb�rer � s�rge for � f� tak i folk som kan gj�re jobben.

\section{Tidsbruk}
Ingen �nsker at du skal brenne lyset i begge ender og ende opp utslitt. Det er fort gjort at studier, foreningsarbeid og fritid blandes litt for mye, og det kan bli slitsomt. Det er derfor viktig � sette grenser. Hvordan du organiserer deg er til syvende og sist kun opp til deg selv.


\chapter{Hovedstyret}
Hovedstyret er Realistforeningens øverste organ mellom Generalforsamlingene.
De planlegger langsiktige mål for foreningen, og har ansvar for at
foreningen som helhet samarbeider godt. Hovedstyret avgjør foreningens
hovedprogram og budsjett, jobber for rekruttering til foreningen, og
arrangerer interne sosiale arrangementer.
% vim:fileencoding=utf-8
% -*- encoding: utf-8 -*-
\section{Formand}

\subsection{Mandat}
Formanden leder Hovedstyret. Formanden er foreningens ansikt utad og skal sørge for at
kontakten med andre studentforeninger, SiO Læringsmiljø og Universitetet opprettholdes.

Formanden har ansvar for å oppdatere denne teksten med nye erfaringer etter fullført verv.
Formanden har også et overordnet ansvar for at alle tekstene i hviteboken blir oppdatert
etter erfaringer og rutineendringer gjennom og spesielt mot slutten av semesteret.


\subsection{Oppgaver}

\subsubsection{Hastesaker ved semesterstart}
\begin{itemize}
\item Brønnøysund. Kan først fikses etter at genforsreferatet er klart.
\item Tilgang til bygg. Liste med navn og kortnummer sendes
      til teknisk-nordre@admin.uio.no. Husk å dele listen i
      to deler: de som trenger tilgang til VB og Abel, og de
      som kun trenger tilgang til VB.
\item Internliste og internkort. Så fort listene med nye funksjonærer er innhentet,
	oppdater internlisten og internkort skrives ut. Bruk en annen farge på internkortene
	enn de tre siste semesterene. Internkort bør være ute senest på semesterstartfesten.
\item Sørg for å få møtevirksomhet i gang, og sørg spesielt for at AS, KS og IS er i rute m.h.p.
	de første arrangementene.
\item Sørg for å bli meldt på de forskjellige maillistene (typisk: kafe, kjellerstyret, as, okonomi, styret).
\item Bytt lås til kontoret. Hvis det ikke ligger flere sett med nøkkelkort i skapet, så bestill nye sett fra TrioVing.
\item Kranførerkurset. Sørg for at dette blir planlagt godt. Husk spesielt på at dette er en introduksjon til foreningen
	og foreningslivet på matnat, og prøv i minst mulig grad å gi inntrykk av at vi skal "kapre" så mange av de nye
	som mulig til foreningsarbeid. Fint om omvisning og øltapping foregår under ordende former, og ikke totalt
	kaos som det pleier å være. Fortell også litt om andre grupperinger av foreningen enn kjellerstyret,
	men ikke for mye om de delene det er uaktuelt å komme rett inn i. Bar, kafe og blæst er hovedmålet,
	mens regi, vaktgruppa og arr godt kan nevnes kort.

	Hold introduksjonstalen kort, og fokuser på fordelene med å være med i en studentforening
	(bredt sosialt miljø, diverse internarrangementer osv) fremfor arbeidsoppgavene.
	Men det er viktig å få frem hvor mye (lite!) jobb det er å være funksjonær.
	Aktiviser alle sammen så mye som mulig. Sørg for at alle som skal presenteres er klar
	over programmet på forhånd.

    Ha internansvarlig med på planleggingen, og sørg for at denne blir tydelig presentert.
    Sørg for å ikke miste listene etter kurset! Sett gjerne en person til å ta
    ansvar for listene og svare på spørsmål når folk kommer for å skrive seg på (feks internansvarlig).
\end{itemize}


\subsubsection{Generelt}
\begin{itemize}
\item Sørg for at Arrstyret, Kjellerstyret og Infostyret holder frekvente møter og at relevante saker blir tatt opp.
	Pass også på at Fagstyret produserer noe, da spesielt at bedriftsansvarlig opprettholder våre kontakter.
	Det aller viktigste er at kommunikasjonen mellom alle disse ledd fungerer, særlig mellom as/fagstyret
	og infostyret. Sørg spesielt for at infostyret starter i god tid med blæsting av store arrangementer,
	slik at det også er tid til litt forsinkelser (som det ofte blir).
\item I hovedstyret bør diskuteres langsiktige saker, eller prekære internsaker. Hold saker som egentlig bør
	omhandles i andre styrer utenfor, og hold orienteringer til et minimum (hvis det ikke eksplisitt berører
	hovedstyremedlemmene).
\item Hold rede på hvilke kontakter foreningen har på fakultetet, teknisk avdeling, eksternt og med andre foreninger.
	Lag gjerne en liste med disse navnene og hvilken tilknytning de har til foreningen.
	Husk å oppdatere denne listen hvert semester. Få hele hovedstyret til å bruke korrespondansepermen,
	slik at vi får samlet viktige papirer og brev på ett sted.

	(Å skaffe samlet liste over foreningskontakter gjennom Cosa Nostra kan ta tid, og vervene skifter fort,
	det kan derfor lønne seg å spørre ledere direkte ansikt til ansikt om kontaktinformasjon).
\item Sørg for at posthyllene blir ryddet ved begynnelsen av hvert semester. Fjern hyller som ikke lenger er i bruk.
\item Innkalling til Hovedstyremøter skal sendes ut minst tre dager i forveien (se lovene).
	Send gjerne først et forslag til dagsorden fem dager før, slik at andre i HS kan foreslå saker de ønsker
	å ta opp. Send deretter innkalling med dagsorden til internelista.
\item Kontorvakter. Sørg for at det på alle hverdager er en kontorvakt i tidsrommet 12-13 (kan tøyes om nødvendig).
	Denne skal hente post (fra postterminalen) og internpost (fra ekspedisjonen på Abel), sortere og notere all
	post i postloggen. Vedkommende skal også rydde kontoret ved nødvendighet (og evt vaske), samt ta imot
	telefonbeskjeder og gi disse beskjedene videre til korrekt mottaker via mail.
\item Pass på at internarrangementer ikke blir glemt, sett helst datoer sammen med resten av programmet.
\item Stille på møter i Cosa Nostra. Kommuniser viktige saker som diskuteres der til Hovedstyret.
\item Sørge for at semesterberetninger er levert inn senest en uke før generalforsamling.
\item Innkalle til Generalforsamling senest 14 dager før generalforsamlingen (se lovene).
	Husk å bestille rom, og kopiere opp dagsorden, semesterberetninger, regnskap og budsjett til alle.
\item Etter generalforsamlingen: Hold på de avtroppende og sørg for at påtroppende i alle verv får nok opplæring.
	Ha klart programutkast for neste semester før eksamensperioden begynner. Etter eksamen kan folk fort forsvinne.
	Spesielt viktig er det å reservere hytte (minst?) et halvt år i forveien. Ved slutten av høstsemesteret, la
	nye og gamle sammen lage budsjett for neste år i samarbeid med ØU. Ved slutten av vårsemesteret, sørg for
	at alle i HS er forberedte på semesterstarten, da den krever sitt.
\end{itemize}


\subsubsection{Andre oppgaver}
I tillegg til de vanlige oppgavene, som noen ganger blir mye og noen ganger lite, har du som Formand mulighet
til å påvirke fokuset i foreningen mot saker du synes er viktig. Noen forslag til oppgaver det neste semesteret:
\begin{itemize}
\item Oppdatere og forbedre hviteboken (skal gjøres hvert semester!)
\item Lage rutiner for bedre opplæring på alle nivåer. Intern "fadderordning"?
\item Brannsikkerhet, rutiner ved brann. Dette er en viktig sak som vi må bli bedre på.
\end{itemize}


\subsection{Tips og triks}

% vim:fileencoding=utf-8
% -*- encoding: utf-8 -*-
\section{Stillingsbeskrivelse for sekretær}

Sekretær er nestleder og Formanns stedfortreder.
Han/hun skal så langt det er mulig assistere og
avlaste Formann i hans/hennes arbeid.

Sekretæren skal i tillegg:
\begin{itemize}
\item Føre referat fra alle Hovedstyremøter og
      sende til Hovedstyret innen påfølgende møte.
\item Sørge for at det til enher tid finnes nødvendig
      kontorutstyr på kontoret og C216.
\end{itemize}

\begin{itemize}
\item Hvert semester oppdatere listen over personer med
      tilgang til C216/Redselskabinettet.
\item Produsere internkort i begynnelsen av semesteret og
      underveis når nye kommer til.
\item Oppdatere internlister dersom internansvarlig mangler.
\end{itemize}


\chapter{Kjellerstyret}
\section{Kjellermester}

Kjellermester er hovedansvarlig for drift av kjelleren. Herunder ligger kafedrift p� dagtid, kjellerkroer og andre arrangementer i kjelleren. 

\subsection{Kjellermesters oppgaver}
\begin{itemize}
\item Bestilling av varer; �l, rusbrus, vin og sprit
\item S�rge for at det alltid er nok av andre n�dvendigheter,
      eks krystallsoda, s�pe til oppvaskmaskin, kluter, s�pe osv.
\item Hold orden p� lagerbeholdningen (til alt utenom kafedrift)
\item Holde orden p� alt som g�r inn og ut av lager
\item F�re Z-rapporter
\item F�re bilag
\item Betale regninger i tide!
\item Holde m�ter i kjellerstyret, og holde seg oppdatert p�
      hva som skjer i de andre styrene (infostyret, arrstyret)
\item Passe p� at de andre gj�r jobben sin.
\item S�rge for at det er nok vekslepenger til en hver tid.
\item V�re med hovedstyret � arrangere kranf�rekurs.
\item V�re med hovedstyret � arrangere internfester.
\item Oppl�ring av neste kjellermester.
\item Oppl�ring av skjenkemestre.
\end{itemize}

\subsection{Under seg har KM}
\begin{itemize}
\item Kaf\'esjef som har med alt det praktiske rundt kafeen � gj�re
      (KM har med andre ord inget reelt ansvar for kaf\'edriften).
\item Skjenkemestre (fire stykker pr i dag). Disse deler ansvaret
      for gjennomf�ring av kjellerkroer, varemottak og organisering
      av Mandager.
\item Utleiesjef. Tar seg av alt det praktiske rundt utleie av kjelleren.
\end{itemize}

\subsection{Hengende ting}
\begin{itemize}
\item Ferdigstilling av kjelleren.
\item Liming av teppefirkanter (rundt baren og i kroken der kj�leskapet stod).
\item Tette sokkelen rundt baren (med silikon).
\item Overbygg over �lslangen.
\item Henge opp s�pedispenser og papirholder.
\item Lage arbeidsinstruks til de forskjellige kjellervervene.
\item Service p� tappeanlegget.
\item Rydding av anti (i samarbeid med kafesjef og arrsjef under ledelse av formann).
\item Riving av toalettene p� Helligste.
\item Kj�ling av Aller Helligste.
\item Rensing av teppet.
\item F� l�s p� alle kj�leskap.
\end{itemize}

\section{Skjenkemester}
Skjenkemestrene har ansvaret for hver sin uke i baren. I en normal uke vil det si Kjellerkro p� fredager, Mandag p� mandager og varemottak p� torsdager.

P� kjellerkroer og andre fredagsarrangementer er det Skjenkemester som har hovedansvaret for serveringen den kvelden. Herunder inng�r: 

\begin{itemize}
\item Ansvar � kalle inn barfunksjon�rer slik at det er nok folk i baren til enhver tid
\item Rydding og renhold i kjelleren f�r under og etter arrangementet (rydde glass, t�rke bord osv.)
\item Holde baren i orden under arrangementet.
\item Fylle p� varer i kj�leskapene.
\item Ta oppgj�ret p� slutten av kvelden
\end{itemize}

Skjenkemestrene har ogs� ansvar for � f�lge opp
oppl�ringen av barfunksjon�rene som ble gitt under Kranf�rerkurset.

I tillegg kan SM bli tildelt oppgaver av Kjellermester s� som:
\begin{itemize}
\item Skrive referat fra Kjellerstyrem�tene.
\item Vaske kluter og mopper.
\item Kj�pe inn s�pe.
\end{itemize}

Samt assistere Kjellermester der det er behov for dette.

\section{Kaf\'esjefens erfaringer}
%\author{Robert Bauck Hamar}
%\date{H�sten 2002---V�ren 2003}

\section{Bestillinger}
Som kaf\'esjef har man ansvar for at kafeen har varer.
\subsection{Tine}
Er v�r viktigste leverand�r. Fra Tine kj�per vi meieriprodukter.
Det er lurt � ha:
\begin{itemize}
	\item Melk. H-Melk til bafler og kaffe. Kj�pes litersvis, 
		men ti liter om gangen.
	\item R�mme. Lettr�mme til bafler. Kj�pes i brett p� seks.
	\item Litago. Kj�pes i ti $\times$ halvvlitere.
	\item Iste og juice. Kj�pes i ti $\times$ halvvlitere.
	\item Norvegia. Kj�pes i femkilospakninger; bemerk dette, idet
		pakninger med to femkilososter har samme varenummer.
	\item Sm�r. Kj�pes i 18 halvkilospakker.
	\item Gudbrandsdalsost.
	\item Kremtopp.
\end{itemize}
Det er lurt � legge merke til at ost og sm�r ikke st�r i bestillingslista.

Jeg har pleid � bruke bestillingsskjemaet som ligger i kaf\'esjefhylla p� kontoret,
men ringt inn bestilling per telefon. Det f�rste for � holde oversikt over hva som
er bestilt. Det andre fordi Tine m� ha bestillinger p� minst 500 kroner for �
levere gebyrfritt. Bestillinger m� v�re Tine i hende senest 10.30 dagen f�r levering.
Leveringer kan man f� p� tirsdager eller torsdager.

N�r det gjelder st�rrelsen p� bestillinger gjelder det � holde seg i nakken. 
Meieriprodukter g�r gjerne fort ut p� dato. Litago holder seg en stund, og
det gj�r ogs� u�pnet ost, men andre ting gj�r det ikke. Det l�nner seg skjelden
� ha mer enn en tipakning juice/iste st�ende. Litago g�r heller ikke altfor fort,
med unntak av sjokolade, som det kan g� en hel pakning av p� en god dag.
Med det nye tellesystemet, som i skrivende stund ikke er helt innf�rt, b�r
det bli enklere � regne ut bestillinger.

Leveringene fra Tine kommer tidlig om morgenen. Som regel blir de dratt inn
av Sandy, en av de ansatte i teknisk avdeling. Ellers m� de dras gjennom
gangen fra Abel til VB. Det gjelder � f� tingene i kj�leskap fortest mulig.
Etter dette skal tralla og melkebrettene settes i vareheisen i kjelleren i
Abel. For � komme dit tar man heisen ned fra underetasjen, g�r inn d�ra
til venstre n�r man g�r ut av heisen, g�r til venstre, rett fram, s� f�rste d�r p�
venstre side.

\subsection{Temperato}
Espressoting f�r man tak av Temperato. Dette er praktisk talt et enmannsforetak,
men den ene gangen jeg bestilte var han kjapp til � levere. Det er greit �
snakke med mannen, og gjerne bestille i store kvanta. Ting man b�r ha:
\begin{itemize}
	\item Espresso. De har litt utvalg. Personlig har jeg foretrukket
		en av de dyrere variantene.
	\item Brunt sukker. Porsjonspakket. F�es i store kvanta.
	\item Mocchasus(?). Til mocchaer. Ta vare p� dispenseren, som
		brukes p� kannene. Kan gjerne \`og pr�ve hvit saus.
	\item Sirup. Se hva som g�r. �nske for eggnogsmak er ytret (St�le).
	\item Glass og kopper.
\end{itemize}
Mannen i Temperato heter Alexander von der Lippe. Han n�s p�: 
\mbox{\textsf{lippe@temperato.no}}, telefon 66~80~87~03, mobil
91~12~28~99 og faks 66~80~36~24.

\subsection{Andre ting}
Det beste er � kj�pe inn det man trenger i store kvanta. L�n bil hos
budsentralen. Pass p� � holde av en bil p� forh�nd. Stedet vi tidligere
har brukt heter Brubekk Storcash, og ligger p� \dots Brubekk. Finn ut 
hvor det er p� forh�nd, og ta gjerne med kart. Det l�nner seg ogs� �
ta med kjellermester, i det minste en annen person. Kundekort er obligatorisk.
P� brubekk kan man kj�pe det meste i store kvanta, kanskje litt billigere.

Egg, br�d og skinke m� n�dvendigvis kj�pes andre steder. Gjerne bunnpris. Det er
enklere for kaf\'esjef om kaf\'epersonalet kan gj�re dette, men det kan
v�re en fordel � gj�re det selv, fordi f�rre trenger � passe p� � huske kvitteringer 
og ha ansvaret for penger. Det blir uansett ting som m� kj�pes p�
Bunnpris til forskjellige tider, s� det er greit at personalet vet hvordan
det funker.

\section{Personalet}
Det er flere kilder til personale. I kafeen har man hatt god tradisjon i � holde p�
personalet gjennom flere semestre. Derfor er det gode grunner til � ta vare p�
personalet. F� folk til � komme p� Mandager$^\text{TM}$, og internarrangementer.
Kranf�rerkurset tror jeg ikke er veldig viktig � nevne mer enn at det l�nner seg
� bli litt kjent med de som er der. En viktig kilde til nye interne har ogs�
v�rt � verve folk selv. Hvis personen en pr�ver � verve er litt interessert i
� bli med b�r det ikke ta lang tid � overtale. Det er \`og lettere � f� med
vennene til folk som allerede jobber i kafeen.

N�r det gjelder Baffeluka er det kaf\'esjefens jobb � skaffe folk. Det aller mest
effektive er � sp�rre folk i kjelleren p� mandager og fredager; det er f� som
svarer p� mail. Bruk et par uker. For din egen del er det enklest � f� kabalen
til � g� opp om folk oppgir flere mulige tidspunkt, men folk flest synes visst
det er best � skrive seg opp ett sted. Foran baffeluka v�ren 2003 brukte jeg
mye tid p� � f� folk til � skrive seg opp flere steder, men ingen m�tte
jobbe mer enn en gang.

Skiftene har tradisjonelt v�rt 10---12, 12---14 og 14---16. Det er dog ingenting
mot � bruke folk annerledes. Det siste semesteret har det v�rt stort press p�
onsdager mellom 12 og 14; det har da v�rt lurt � ha tre personer p� 12---14-vakta.
Det kan man regne med � m�tte gj�re framover ogs�.
En lur ting er � sette opp et skjema tidlig slik at alle vet n�r man skal jobbe.

\subsection{Organisering}
Det er flere ting � organisere start og slutt. En l�sning som har fungert det
siste semesteret er at en av de som st�r klokka 10 har ansvar for � �pne.
For � overlevere n�kkelen b�r denne personen ha tilgang til kontoret.
RF tenker p� � kj�pe et l�sbart n�kkelskap, og dette kan brukes. Inntil
da f�r man finne p� noe annet lurt.

N�r det gjelder � stenge har jeg valgt � gj�re det selv. En grunn til det
er at f�rrest mulig skal trenge � kunne gj�re feil i regnskapet. Men
det g�r selvf�lgelig an � delegere ansvar videre. Jeg vil bare peke
p� at skjenkemestrene ogs� tar oppgj�rene sine selv.

\section{Diverse}
\subsection{Baffeluka}
For � f� Baffeluka til � bli ei god Baffeluke, er det noen ting
som m� p� plass. Jeg har allerede nevnt personale. Videre er det viktig
at folk vet det er Baffeluke. Bl�st kan st� for bl�sting. Send gjerne en 
mail til bl�st eller bl�stsjef. 

Det aller viktigste er at Regi vet n�r
Baffeluka er. Jeg snakker av erfaring: Send mail, ikke bare snakk med folk.
Gjerne ogs� p�minnelser. Som Sturle sa: �Regi er treige,� de f�lger ikke like
godt med som andre grupperinger. 

Nummer tre: filmer. Lene Kittelsen er gjerne en person som er kyndig til �
plukke ut filmer. Finn filmer i god tid, slik at det er lettere � f� tak i dem.
De interne kan gjerne ogs� bli h�rt. Det holder med to filmer per dag, eventuelt
\'en dersom den er lang. En god baffelukefilm er en film man kan se p� uten � se
hele. Den siste filmen er alltid �Kellys helter.� Regi har den p� VHS.

Baffeluka grunner i at kjelleren var �pen
fordi folk ville se OL. OL og Ski-VM pleier gjerne � g� sammen
med Baffeluka p� v�rsemesteret. P� h�stsemesteret kan man gjerne ogs� vise
TV med et idrettsarrangement.

\subsection{Ting som b�r gj�res}
Flere ting b�r gj�res. F�rst og fremst b�r man lage en stillingsinstruks for
personalet. Den b�r for eksempel inneholde instruksjoner om hvordan maskiner
fungerer og andre ting. Gjerne ogs� Baffelr�reoppskriften. Etter et 
kjellerstyrevedtak er oppskriften hemmelig. Som en morsom avledningsman�ver,
b�r man kanskje fake oppskriften p� en m�te alle personale kjenner til.
For eksempel kan �let byttes ut med Bache Gabrielsen VSOP.

Man b�r ogs� holde orden i skapene. En id\'e er � lage oversikt over hva som 
befinner seg i skapene. Ting som andre kan ha glede av, for eksempel vannkoker
og kaffe, b�r ikke l�ses inn. Apropos kaffe: Jeg har innf�rt en ordning der
folk som tar kaffe putter penger p� et glass. Glasset kan t�mmes av og til,
og pengene brukes til ny kaffe.\footnote{Denne kaffen er allerede inne i regnskapet,
s� man trenger da ikke kvittering/bilag.} Pass p� at folk ser glasset, og forst�r
ordningen.

N�r det gjelder kj�leskapene, har man lite plass. Hvis det trengs, kan man ta
ting som ikke trenger � holdes kaldt ut. For telling og praktiske form�l, er
det lurt � holde orden.

F� tak i noe til oppvaskmaskina, som kan h�ndtere bestikk. Jeg har sett hele 
rister med finmasket nett. Til v�rt forbruk holder det nok med en bestikkurv
fra en vanlig oppvaskmaskin.

\subsection{Praktisk info}
Oppvaskmaskinen er n� koblet opp mot s�petanken, slik at s�pen egenhendig
blir sugd inn i maskina. Den m� nok kalibreres, og det st�r en guide til
dette i manualen. P�se at det er nok s�pe i dunken. S�pe kj�pes ogs�
p� Brobekk. �lvaskeren bruker fortsatt krystallsoda.


\chapter{Økonomiutvalget}
% vim:fileencoding=utf-8
% -*- encoding: utf-8 -*-
\section{Økonomiutvalget}

\subsection{ØU-leder}

\subsubsection{Stillingsbeskrivelse}

\begin{itemize}
\item{} leder økonomiutvalget og derfor overordnet ansvar for
  utvalget. Dette vil si å kalle inn til møter, finne møtedager som
  passer for de fleste.
\item{} skal ha oversikt over RFs totale økonomi. Dette betyr at
  vedkommende skal kunne gjøre rede for RFs økonomiske situasjon enn
  når dette måtte forekomme.
\item{} skal sørge for at de forskjellige utvalgene fører en
  ansvarsfull økonomi. Dette gjøres gjerne ved å lære utvalgene opp
  etter RFs gjeldende regler, rutiner og system dersom det er behov
  for dette.
\item{} informere de andre utvalgene hvordan de ligger ann ifølge
  budsjett. Dette er viktig for at man til enhver tid holder seg
  akkurat eller i nærheten av budsjettet.
\item{} holder bilagskurs for de som gjennom sitt verv trenger å
  skrive bilag. Dette er en meget viktig arbeidsoppgave. Grunnen til
  dette er at man ønsker en så homogen som mulig bilagsskriving slik
  at regnskapet blir oversiktlig å lese for eventuelle utenforstående.
\item{} delegerer oppgaver til medlemmene. Som en god leder gjelder
  denne oppgaven som alltid. Det er viktig å prøve å motivere
  medlemmer til å synes det er greit å være med i ØU.
\item{} passer på at regninger blir betalt. Alle regninger skal
  betaler før de går ut på dato, rimelig nok.
\item{} sørger for at MVA blir betalt. Dette er et veldig viktig
  punkt. MVA-oppgaven/omsetningsoppgaven er RF sitt ansikt utad mot
  skattemyndighetene. Vi får igjen skatt på de varene iv kjøper inn som er
  momsbelagt og vi betaler moms for de varene vi selger igjen. I
  RF-sammenheng gjelder dette mye leskende drikker.
\item{} snakke med bank og andre instanser i spørsmål som har med
  økonomi å gjøre. Siden de som sitter i ØU sjelden har så mye
  erfaring med økonomi, er det viktig å være etterrettelig å prøve å
  få ut informasjon der det er mulig. Mest av alt gjelder det å spørre
  tidligere ØU-medlemmer (derav ledere og sekretærer). 
\item{} skrive ut bilagsforsider. Dette er en ren
  rutinejobb. Bilagsforsidene trengs for de fleste typer bilag.
\end{itemize}

\subsubsection{Rutinebeskrivelse}

\begin{itemize}

\item{} Etter at det nye hovedstyret er satt, har ØU-leder ansvar for å kalle inn de
  relevante og gå i banken for å fornye tilgang til kontoene. Dette
  kan bare gjøres med et underskrevet generalforsamlingsreferat i
  hånden. ØU-leder kaller sammen de i hovedstyret det gjelder, for
  øyeblikket Kjellermester, Arrmester og Formann. ØU-leder skal ha
  tilgang til alle kontoer. Kjellermester skal ha tilgang til
  kjellerkonto, Arrmester arrkonto og Formann sparekonto. Sparekontoen
  skal bare kunne manipuleres med både Formann og ØU-leder tilstede.

\item{} \textbf{Betaling av fakturaer} Dette skjer ved hjelp av
  postgiro og signeres på vanlig måte når dette er ordnet i banken.

%\end{itemize}
\item{} Innlevering av MVA-/omsetningsoppgave. Skattemyndighetene skal ha kontroll
  over all omsetning RF har. Derfor leverer vi MVA-oppgave hver annen
  måned. Alle ark som trengs blir sendt fra Flykesskattesjefen og
  innleveringsfristen er som regel den 10. i hver partallsmåned. I
  forbindelse med innleveringen er det innført interne
  kontrollrutiner. Disse er å finne på et ark i en av
  permene. De viktigste rutinene er
  \begin{itemize}
  \item{}telling av alle kasser. For å finne ut om summen i kassen
    stemmer med regnskapet.
  \item{}bank avstemning. Se at det vi har ført i bankpostene faktisk
    stemmer med bevegelse i banken.
  \item{}leverandør- og kundeavstemning. Dette betyr at man må gå
    gjennom alle kunder og leverandører (regnskapet) og sjekke om de går i null
    eller ikke. Dette er også noe som skal gjøres etter hver
    postering.
  \item{}sjekk om penger fra kassa er ført inn i regnskapet.
  \end{itemize}
\end{itemize}
Denne jobben tar faktisk litt tid . Det er derfor viktig
å begynne med disse sjekkrutinene i god tid før
MVA-oppgaveinnlevering.
\begin{itemize}
\item{}\textbf{Avslutting av perioder} Hver periode (måned) skal
  avsluttes i regnskapsprogrammet. Dette skjer som regel i forbindelse
  med momsinnbetalingen. 

\item{} Avslutting av år. Rutinene er de samme som for avslutning av
  perioder. Den eneste store forskjellen er at man skal sjekke om hele
  året stemmer. Eventuelle differanser, som av forskjellige grunner
  ikke kan føres bort, bør beskrives i et vedlegg til
  regnskapet. Permer med samtlige bilag settes opp i bokhylle for
  eventuelt ettersyn.

\item{} Putte bilag i perm. Alle bilag puttes i en av 3 permer før
  føring. Hver av permene har sin egen nummerserie som identifiserer
  typen bilag ført.

  \begin{itemize}
  \item{} \textbf{hovedboksbilag} Alle bilag som har med bevegelse i
    bank å gjøre. Dette gjelder f. eks. nattsafeposetellinger levert fra
    bank via fax og kvitteringer for betaling av giroer osv.
    
  \item{} \textbf{kontantsalgbilag} Alle bilag som har med salg i
    baren/kafeen å gjøre. Disse bilagene inneholder hvor mye penger som
    ble talt opp kontra hva kassen mente. Differansen blir
    svinn. Pengene som gikk inn i kassen føres mot riktige salg-poster i
    regnskapet. Kvittering for nattsafepose skal stiftes på bilaget,
    samt Z-rapporter fra kassa. Summen som var opptalt i kassa skal være
    lik summen på nattsafeposen. Dersom dette ikke er tilfelle skal den
    som skrev bilaget kontaktes for å gi en forklaring på forskjellen.
    
  \item{} \textbf{inngående-faktura-bilag} Andre bilag;
    betalingsmeldinger, utlegg, osv. Disse føres mot leverandører og
    kunder. 
    \begin{quote}
      \textbf{Når en leverandør har levert varer til Kjellern skal kjellermester betale gjeldene faktura.
        Kopien (merket kopi eller kopiert av RF) skal gis
        ØU-leder. Dette for å unngå dobbeltføringer.} 
    \end{quote}
  \end{itemize}
\item{}Ved semesterstart skal ØU-leder skifte alle koder på alle safer
  for å hindre at de som ikke skal ha tilgang til safer har det.  
\end{itemize}

\subsection{ØU-sekretær}

\subsubsection{Stillingsbeskrivelse}

\begin{itemize}
\item{} Skriver referat fra ØU-møter hvis behov. Som regel er det ikke
  behov for referater. De fleste ØU-møter går ut på å føre bilag.
\item{} Hjelper til med føring av bilag og hjelper
  til med å lære opp nye medlemmer. Det å styre et regnskap tar litt
  tid å sette seg inn i. Derfor er det veldig viktig at nye medlemmer
  får god opplæring.
\item{} Sørger for og administrerer full lagertelling. For å ha
  kontroll over svinn i lageret er det viktig med månedlig
  lagertelling. Å administrere telling innebærer å avtale et passende
  tidspunkt med kjellermester samt utføre tellingen. Resultatet av
  tellingen føres på eget ark (finnes i exel-format), hvor all
  summering gjøres automatisk. Tellingen føres også inn i regnskapet
  for å kunne sammenligne poster for innkjøp og salg av diverse
  varer. Lagertellingsresultatet skal typisk være differansen mellom
  innkjøp og salg.
\end{itemize}

\subsection{Rutinebeskrivelse}
\begin{itemize}
\item{} møte opp på møter eller melde forfall.
\end{itemize}


\subsection{ØU-medlem}
\subsubsection{Stillingsbeskrivelse}

\begin{itemize}
\item{} lærer å føre bilag og å forstå hvordan rutinene er
\item{} hjelper til 
\end{itemize}

\subsubsection{Rutinebeskrivelse}
\begin{itemize}
\item{} møte opp på møter eller melde forfall
\end{itemize}

\subsection{Erfaringsprotokoll}
Erfaringsmessig er det vanskelig å følge alle rutiner til punkt og
prikke. Siden dette er en studentforening hvor vervene faktisk gjøres
på fritiden og utenfor skole, kan man ikke forvente at alt skal
gjøres. Stillings- og ruteniebeskrivelsen er ment som veiledende. Det
viktigste er at ØU har kontroll over alle bevegelser i regnskapet og
en oversikt over foreningens økonomi, samt leverer inn omsetningsoppgave. 


\chapter{Arrstyret}
% vim:fileencoding=utf-8
% -*- encoding: utf-8 -*-
\section{Arrsjef}

Realistforeningens (RF) arrsjef har ansvaret for forberedelse og avvikling
av RF's arrangementer.

Arrsjefen leder arrstyret, som i tillegg består av
regiformann, vaktgruppesjef, blæstsjef og arrmesterene. Arrsjefen
skal koordinere denne gruppens arbeid på en slik måte at alle vet hvilke
arrangementer som er forestående, og hva det er forventet at de skal gjøre
i sammenheng med disse.

Arrstyret bestemmer hvilke arrangemter RF skal avholde, og er ansvarlige
for avvikling av disse, f.eks.
kjellerkroer, temafester, konserter o.l., Høstpølse, Psildefrokost og
julebord. I tillegg er arrstyret ansvarlige for Danmarkstur og hyttetur.
Internfester er, med mindre arrsjef i samarbeid med hovedstyret har blitt
enige om det, ikke arrstyrets ansvar.

Under planleggingen av et arrangement skal arrsjefen holde kontakten
mellom de forskjellige grupperingene i arrstyret. Siden arrstyret er
stort og driver med svært forskjellige ting kan dette ofte vise seg
vanskelig, og arrsjefen må derfor være fleksibel i sitt arbeide.

Arrsjefen har ansvaret for å holde kontroll på RF's medlemskapssalg. Dette
innebærer å sørge for at det blir solgt medlemskap på ett arrangement, at
medlemslisten til enhver tid er ajourført, og å sørge for at det finnes
medlemskort.

På ett arrangement er det arrsjefens ansvar å sørge for at det personalet
som kreves for å drive arr. er tilstede. Dette innebærer at det skal være
vakter tilstede (hvis nødvendig), personer til å håndtere lys og lyd (hvis
nødvendig, mannskap til å betjene baren (hvis nødvendig), samt en
hovedansvarlig. Arrsjefens jobb må her anses som gjort hvis de ansvarlige
for disse undergruppene har fått tilstrekkelig instruks i god tid.
Arrsjefen skal ikke lastes for andre grupperingers manglende evne til å
gjøre sine ting.

Før et arrangement skal det også blæstes. Det er arrsjefens jobb å
informere blæstsjefen om hvilke arr som er nært forestående. Blæstsjefen
skal da ha tilstrekkelig informasjon om arr til at han/hun kan blæste arr
skikkelig.

Arrsjefen skal velge ut arrmestere til sin arrgruppe. Disse innstiller han
så for hovedstyret som så stemmer.


\section{Arrmestere}

Arrmestere blir utpekt av arrsjef og innstilt for hovedstyret.
Arrmesterene sitter i arrstyret.

Arrmestere skal assistere arrsjefens i hans oppgaver.

De skal assistere i
planleggingen og forberedelsene til et arrangement, samt hjelpe til under
arrangementer. Arrmesteren skal også jobbe som hovedansvarlig (se egen
instruks) på RF's arrangementer. Arrmestere må derfor ha nødvendig
kunnskap om brannsikkerhet.

\section{Tradisjonelle arrangementer}
Her følger en rekke arrangementer som RF har mer eller mindre lange
tradisjoner på å holde.

Dette må noen skrive.
%\input{kildekode/arr/harryaften}
%\input{kildekode/arr/psildefrokost}
%\input{kildekode/arr/lasvegas}
%\input{kildekode/arr/høstpølsefest}
%\input{kildekode/arr/cowboyaften}
%\input{kildekode/arr/julebord}

\section{Erfaringsprotokoller}
Denne seksjonen kan med fordel fylles på med flere erfaringer og tips
i forbindelse med arrangementer.
% vim:fileencoding=utf-8
% -*- encoding: utf-8 -*-
\subsection{Karneval 14/02/2003}

\subsubsection{Innkjøp}
5*4meter bannerstoff av ymse farger kr 960.-
hatter og fløyter og girlander og premier kr 140.-

\subsubsection{Gjennomføring}
Banner ble hengt opp en uke før, plakater en til to dager før
Vi kjøpte inn "inngrediensene" torsdag kveld. Fredag fra ca klokka
16.00 begynnte vi å pynte i  kjelleren. Vi hang opp stoffer og
ballonger langs vegger og  stolper og girlander i baren. 
Gjestene kom til normal tid 18.59, men de ferreste utkledd (stort sett
interne som var utkledd).  Ellers var det senere på kvelden kåring av
beste kvindelig og  mannlige kostyme. 

\subsubsection{Bemerkninger}
Det ble litt stress med pyntinga på slutten, men det
ordnett seg til slutt. Hvis det skal være lys på karneval neste gang,
må man ta kontakt med regi før ett par dager før. De har ett lysnett
på lager, som ble brukt, men ønskes noe større må det bestilles. For å
henge opp stoffene ble det brukt tråd og tau til
forsterkning. Blæstinga må skje mye før, og gjerne med plakater. Det
kom flest interne som hadde kledd seg ut, for å ordne på dette, blæst
før!

% vim:fileencoding=utf-8
% -*- encoding: utf-8 -*-
\subsection{Alle Festers Mor 14. mars 2003}

Det ble kjøpt inn t-skjorter. 
Vi hadde et sammarbeid med FF og Studentparlmentet på betong. Vi var i
hovedbaren. En vanskelig bar å pynte, men det virket som det kom en
del folk dit. Pynting var ikke mer enn et banner som det sto
bjørnehiet på. Ellers lagde vi noen flyers med info om RF og hvordan
bli medlem. Barpersonalet brukte de nye t-skjortene.
Vi hadde sammen med FF ansvaret for lyd og lys. Regi var her veldig
behjelpelig, og tok seg av størsteparten her. Vi fikk leid inn en
kjempe discokul og mengder av lys effekter. Vi leide også inn ekstra
lyd, da de ikke har nok på betong. Det var to dj-er som spilte fra
oss, og en fra FF.
Arrangementet gikk bra, men det er en god ide å begynne planleggingen
mye før, og ta kontakt med regi så tidlig som mulig.

Arrsjef \\
Olav


\chapter{Fagstyret}
% vim:fileencoding=utf-8
% -*- encoding: utf-8 -*-
%Denne fil vil trenge en
%\usepackage{amsmath,url}
%i preamble.
%
% \chapter{Fagstyret}
Om fagstyrets utforming og mandat:
\begin{quote}

    % Ajour per H10
    {\Large § 11 FAGSTYRET}
    
    \begin{itemize}
        \item[a)] Fagstyret ledes av Fagsjef, og har følgende andre medlemmer:
        Populærvitenskapelig sjef og Bedriftskontakt.
        
        \item[b)] Fagstyret har ansvar for foreningens faglige profil og
        kontakt med næringslivet.
        
        \item[c)] Fagstyret skal fremlegge forslag til foreningens faglige
        program for Hovedstyret.
        
        \item[d)] Fagstyret har ansvaret for å rapportere regnskapsrelevant
        informasjon til Økonomiutvalget og overholde vedtatte budsjetter.
    \end{itemize}

\end{quote}

\section{Fagstyrets oppbygning}
\subsection{Fagsjef}
Fagsjef er sjef for fagstyret, og sitter i tillegg i Hovedstyret.
Ved siden av å lede fagstyret er
fagsjefen som regel også enten Pop.Vit.-ansvarlig eller Bedriftskontakt.
I motsetning til de fleste andre verv er fagsjef en ganske fri stilling; Ting
som gjøres skjer oftest på eget initiativ.

\subsection{Bedriftskontakt}
Bedriftskontakten har ansvar for å pleie kontakten mellom arbeidsgivere og
foreningen, bl.a.~ved å arrangere bedriftspresentasjoner og ekskursjoner.

\subsection{Populærvitenskap-ansvarlig}
Pop.Vit.-ansvarlig har hovedansvar for å skaffe foredragsholdere og debattanter,
og å finne relevante temaer.

\subsection{Panikkansvarlig}
Panikkansvarlig har ansvar for å arrangere panikkhjelp.

\section{Ansvarsområder}
\subsection{Foreningens faglige profil}
\begin{quote}
    % Allerede de gamle grekere?
	Allerede de aller tidligste lover slo fast at foreningens formål er delt: en
	faglig og en sosial del. De tidligste møter i foreningen ble innledet
	med foredrag og debatter før de gikk over i fest og moro. Selv om den
	faglige delen i dag står noe svakere, er den en viktig del av
	foreningen.
\end{quote}

For foreningens faglige profil arrangeres det fortsatt foredrag, kurs og andre
ting.

For å arrangere noe kreves følgende:
\begin{itemize}
	\item Personer, for eksempel foredragsholdere. Dette er den mest
	krevende faktoren.
	\item Lokale. 
          Kjelleren er et fint lokale, den inneholder alt du ønsker av
          mikrofoner, prosjektorer etc. Ta kontakt med utlånsansvarlig i KS.
          Ellers kan Universitetets lokaler reserveres vha. mail.
    % XXX Komstyret: Denne må fikses hvis kommunikasjonsstyret innføres
	\item Blæsting. Det er viktig å gi blæst informasjon, de fikser også event
          på \texttt{facebook}. Også webansvarlig bør få informasjon til
          hjemmesiden.
	\item Gave til foredragsholder, for eksempel en flaske vin. 
          Taes ut fra kjelleren; før bilag og \emph{lever lapp} til
          Kjellermester. Se også \ref{fagstyret:sec:bilag}.
	\item Mandag$^\text{TM}$ kan gjerne gjennomføres etterpå, hvis
          arrangementets kjer på kveldstid.
\end{itemize}

Populærvitenskapelige foredrag er typisk det vi arrangerer oftest.
Man bør strebe etter å arrangere noe et par ganger per måned.
Det er viktig å ha noenlunde oversikt over hva som skjer av foredrag ellers
på fakultetet, så man ikke arrangerer ting samtidig som andre.

Det finnes mange potensielle foredragsholdere på universitetet, det gjelder bare
å lete.
Statlige tilsyn og foretak er også glade i å komme for å fortelle om det de
driver med.
% Hjelp til å lete kan finnes i permen på kontoret som er merket
% ``Fagstyret''. Der har vi informasjon tilgjengelig om hvem som har holdt 
% foredrag tidligere, samt kontaktinformasjon. For 2003....
I tillegg har vi en avtale med Tekna, som gir oss
hjelp til å finne kompetente folk som kan et visst tema. 
Det kan være et triks å velge ``dagsaktuelle'' temaer:
Ting som blir skrevet om i avisen osv. ---
debatter trekker alltid mange mennesker.

NB: det er \emph{helt sentralt} å ha en god \emph{ordstyrer} når man
arrangerer en debatt, ellers har det en tendens til å skli ut.

\subsection{REAL frokost}
Den vanligste pop-viten er i skrivende stund i form av REAL frokost:
Fagstyret (evt. forsterket av Arr) står opp tidlig (de er på blindern
ca.~0730) og lager frokost, som så serveres til alle som ønsker, samtidig som
det foregår en debatt eller et foredrag.
Veldig populært!
Se også \ref{fagstyret:sec:REALfrokost-oppskrift}.

\subsection{Panikkhjelp}
Fagstyret arrangerer for tiden panikkhjelp i første\-års\-emner på matematisk:
MAT1001, MAT1100, MAT-INF1100 og STK1000 i høstsemesteret, og MAT1110, MAT1001
 og STK1100 i vårsemesteret. ``Panikkhjelperene'' er som regel gruppelærere i
fagene, og personer i foreningen som har hatt fagene før.
Panikkhjelpansvarlig er ansvarlig for å gjennomføre panikkhjelp.

Matematisk institutt pleier å sponse panikkhjelpen med penger til mat til
``panikkhjelperene''.

Tidligere har det også vært arrangert turer til forskjellige steder.
Disse kan med fordel arrangeres i samarbeid med relevante arbeidsgivere, 
men man kan også dra på utstillinger,
eller f.eks.~``Senkveld'' på Teknisk Museum.

\subsection{Kontakt med næringslivet}
Å ha kontakter i næringslivet er viktig. Foreningen kan for eksempel tjene
penger på å holde bedriftspresentasjoner i kjelleren. I tillegg kan man gjøre
det mer fordelaktig for jobbsøkere å ha RF på CVen.

RF har en avtale med Tekna. I den forbindelse, har vi fått en liste med bedrifter
der Tekna er representert. Det kan da typisk være lurt å holde kontakt med
Tekna-tillitsvalgt i disse bedriftene, slik at denne personen kan hjelpe oss
med å overtale bedriften til å komme på besøk.

Det finnes mange gamle interne som nå jobber i relevante bedrifter, de kan være
nyttige når man skal ordne bedriftspresentasjoner og lignende.

% Dessverre står det ikke så mye her enda. Mye jobbing står igjen for å bygge opp
% denne delen. Nå skal også Biørnegildet 2004 arrangeres, og jeg håper at
% sponsoransvarlig Lars Warholm kan være behjelpelig med å skrive mer under
% denne subsubsection.

\section{Om sosialt}
Et av Realistforeningens viktigste formål er å bygge et sosialt miljø
mellom forskjellige realister, og det er derfor også et formål for
fagstyret; Det er viktig at det er godt samhold internt i fagstyret, og
at fagstyret er en integrert del av RF.
Alle nye fagstyremedlemmer bør oppfordres til å komme på internfester og Mandager.
I tillegg bør det arrangeres et par bare-fag sosiale arrangementer i løpet av
semesteret, som for eksempel fagstyremiddag og fag-på-fyll.
Fagstyremiddag kan forøvrig med fordel kombineres med det å skrive semesterberetning.

\section{Tips og triks}
\subsection{Fagstyrespesifikt om bilag}
\label{fagstyret:sec:bilag}
Bilag føres \emph{hver gang} penger skal innom foreningen. 
Fagstyret har fortiden følgende kontoer:
``7120 Pop.~Vit.~og fagstyreavgifter'',
``7062 REAL frokost'', og
``7073 Gave til foredragsholder''.
Sistnevnte er ment for  betale bl.a.~vin til foredragsholder,
den midterste bør være relativt åpenbar,
og den førstnevnte brukes til alt annet.
I tillegg finnes kontoen
``7060 Eksterne arrangementer'', men den er skal egentlig bare brukes av Arr.
Hvis du ikke vet hvor du finner bilagsskjemaer, kan du spørre ØU-sjef.
(Vedkommende vil mest sannsynlig også fortelle deg hvordan man fører bilag.)

\subsection{Blæsting}
For at et arrangement skal bli en suksess, er det nødvendig at det kommer
folk. For at det skal komme folk, er det nødvendig at de informeres om
arrangementet.
Blæst er eksperter på dette. De både lager og henger opp plakater, og lager
dessuten event på facebook osv.
Det er viktig å alltid holde blæst informert om kommende arrangementer.
For tiden løses dette bl.a.~med at blæstsjef er med på fagstyre-listen.
Det er allikevel viktig med direkte kommuniksajon mellom Fagsjef (eller
den som er ansvarlig for et spesifikt arrangement) og Blæstsjef.

Det kan alltid være lurt å snakke med foredragsholderen for å finne ut
hva som skal skje. Han/hun kan skrive et kort sammendrag, og kanskje komme
med bildemateriale som kan brukes på plakater/banner og nett.

Noen former for blæsting kan fagstyret stå for selv. Det er lurt å få 
arrangementet ``på plakaten'' i Universitas. Dette gjøres ved å sende
en mail til \url{til-plakaten@universitas.no}.
Det eksisterer en liste som heter \url{informasjon@rf.uio.no}
Sender man informasjon til denne lista, blir
det sendt til alle interne, samt endel eksterne som ønsker informasjon.
Listen er moderert, så enkelte ting kan ta lang tid.

\subsection{Oppskrift på et faglig arrangement}
\label{fagstyret:sec:REALfrokost-oppskrift}
Når en REAL-frokost-arrangør skal arrangere en REAL-frokost, er det en del
ting som må gjøres:
\begin{itemize}
    \item Skaffe foredragsholder eller debattanter
    \item Skaffe lokale
    \item Kontakte blæst, fikse blæsting
    \item Kontakte Regi, fikse Lyd/Lys/Etc.
    \item For REAL frokost: Handle
    \item For REAL frokost: kontakte teknisk avdeling og si at man skal
          lage mat, så de ikke dukker opp uventet og er ufine.
    \item Utføre
    \item Føre bilag (husk å be om kvitteringer når du kjøper noe)
\end{itemize}


\chapter{Infostyret}
\section{Stillingsbeskrivelse for infosjef og bl�stgruppa}

Infosjefs hovedoppgave er � informere eksterne og interne om hva som skjer i kjellern.
Dette er en viktig oppgave, som ikke kan gj�res av en person alene.
Derfor har infosjef en webansvarlig, en listeansvarlig og en bl�stgruppe med seg p� laget.
Bl�stgruppa burde v�re p� minst ca 5 personer, jo flere jo bedre. Medlemmer b�r helst rekruteres p� slutten av forrige semester, men semesterstart kan ogs� v�re en fin tid for � utvide gruppa. Spesielt gjelder dette h�stsemesteret, n�r det er mange nye studenter. Det kan v�re greit � bruke en del tid p� � snakke med nye studenter p� de forskjellige arrangementene, og informere om RF generelt, og bl�st spesielt.

Infosjef kaller inn til bl�stm�ter n�r det trengs, og helst i god tid f�r m�tet
(min. 2 arbeidsdager f�r). M�tene brukes som oftest til � fordele oppgaver,
eller diskutere forskjellige virkemidler man vil bruke i tiden fremover.
Bl�sting b�r v�re oppe senest en uke f�r arrangementet, s� det gjelder � planlegge
tidsbruken god tid i forveien.


Vanlige PR-metoder:
\begin{itemize}
\item Banner p� Frederikke
\item Plakater over hele Blindern
\item Flyers
\item Stuntbl�sting
\item RF's mailingliste
\item Nettsiden v�r
\item Universitas
\item UiOs nettsider
\end{itemize}


\subsection{Bannere}
Husk � kapre plass s� fort som mulig tidlig i semesteret. Det er to muligheter for � lage bannere.
\subsubsection{Male selv}
(Dette er den tradisjonelle metoden for � lage bannere)
N�r man skal male banner kan det v�re lurt � tegne opp med kritt f�rst. Da 
skriver man ut et lysark (som ligger i bl�stskapet) med tekst og bilde p�.
For s� � henge opp bannerstoffet og rette overheaden mot det, s� er det bare
� tegne av. Deretter legger man f�rst ut plast p� dansegulvet, for s� � male
banneret opp� plasten. Plasten ligger typisk i bl�stskapet, eller i en av 
benkene i kjeller'n.
\subsubsection{Bestille}
Fler og fler begynner � bestille trykte bannere p� plast. Dette ser proffere ut, og kan v�re mer i�ynefallende. Det er dog noe dyrere, og b�r ikke brukes for ofte. Kan f.eks. brukes for hovedbanneret som skal henge der hele semesteret.
\subsubsection{Henge opp banner}
For � henge opp banner p� Frederikke kan det v�re kjekt med en lang 
gardintrapp. Det kan man enten f� l�nt i resepsjonen p� idret, eller i
informasjonsdisken p� Abel. Stiftemaskin ligger i bl�stskapet p� VB.


\subsection{Plakater}
For � f� til fine plakater er det lurt � f� tak i noen som enten
er flinke til � tegne for h�nd eller er dyktige med datagrafikk.
Gjerne begge deler. Det kan kanskje v�re en fordel � ha et gjennomg�ende
tema i plakatenes utseende, men ikke lag dem for like. Da blir
de fort oversett etter hvert.

For � f� skrevet ut plakater kan printeren ''Goliat'' benyttes. Den st�r p� 
ifi, og man m� fylle ut et bestillingsskjema for � bruke den. 
Bestillingsskjema kan man finne p�:
\url{http://www.usit.uio.no/it/utskrift/bestilling\_goliat.html}
annen nyttig informasjon om goliat kan finnes p�:
\url{http://www.usit.uio.no/it/utskrift/goliat.html}
eller ring operat�rene p� (228)52501.
Det finnes ogs� printere for � skrive ut i farger, og for � skrive ut p� begge sider p� ifi.

Bl�stfunksjon�rene b�r tildeles ruter p� Blindern der de har ansvaret
for at det henges opp plakater. Infosjef kan f.eks legge bunker med
plakater p� kontoret eller i glassmonteren utenfor kjellern, og s�
kan funksjon�rene henge dem opp n�r hver og en har tid.

Husk � ta ned plakater etter at arrangementet er ferdig. Dette er
veldig viktig, men blir sjelden gjort. Det ser uprofesjonelt ut n�r det
henger oppe gamle plakater, samt at folk kanskje ikke gidder � se p� dem
fordi de som regel er gamle.

Plakater b�r henges opp i god tid f�r arrangementet! 1-2 uker kan v�re passe, men hvis de henger for lenge blir de klistret over innen arrangementet begynner. Is�fall kan man g� en liten ekstratur noen dager f�r og henge opp ekstra plakater.

\subsection{Flyers}
N�r man skal lage flyers er det lurt � bruke kuttemaskin for � dele opp
arkene man har skrevet dem ut p�. Det finnes en kuttemaskin p� ifi,
og det skal ogs� skal finnes en hos SiO L�ringsmilj�.

I forkant av store arrangementer g�r det an � legge ut l�pesedler p�
lesesaler. Man m� da sp�rre lesesalsvaktene (evt. RSU hvis man vil legge
ut i VB) og man m� komme n�r de �pner. Man m� ogs� samle dem inn igjen n�r
lesesalen stenger. Regner med at det samme gjelder for termstuer.

Andre steder der flyers kan deles ut er p� spisestedene p� Frederikke, og i andre fellesarealer der folk oppholder seg (slik som i aulaen p� VB).

\subsection{Web og mail}
Nettsidene v�re er avhengige av oppdatert informasjon for at de skal v�re
interessante. Hvis det hele tiden ligger aktuell informasjon der, kommer
folk igjen senere. Lag derfor sm� tekster til alle arrangementer, og
send dem til webansvarlig sammen med, eller i etterkant av, programmet
for semesteret. Nettsidene v�re viser de neste to-tre ukene med program
automatisk, s� hele semesterets program burde legges inn med en gang det
er klart.

Send ut informasjon p� informasjon@rf.uio.no. Denne lista b�r brukes i
forkant av alle slags arrangementer. En jevnlig nyhetsmail (f eks ukentlig)
er en god ide. Det er ogs� viktig � v�re aktiv med � f� interesserte til � melde seg p�.
Lister p� oppslagstavler (husk � hente dem inn!), flyers med nettadressen v�r,
lister i d�ra ved arrangementer, lister i kafeen, mulighetene er mange for
� finne interesserte.

Legg ut informasjon p� UiO's "oppslagstavle". http://www.uio.no/tavle/
Her m� man f� tilgang til � legge ut informasjon, og det f�s ved � sende
mail til de som drifter tavla. Info finnes p� websiden.


\subsection{Eksterne informasjonskilder}
F� med alle arrangementer (kaf�, foredrag, kveldsarr) i "P� plakaten" i
Universitas, det er gratis, og det mest selvf�lgelige stedet av dem alle
� v�re synlig. Straks programmet for semesteret er ferdig, kan det v�re lurt
� sende inn hele lista med korte tekster (liten plass) til hvert arrangement,
slik at de kan legge det inn i systemet sitt.

Ha gjerne en spalte i Universitas i forbindelse med store arrangementer.
Dette koster penger (ca 800,-). Informasjon om dette st�r i Universitas.

Ved store arrangementer hvor foajeen er i bruk kan det v�re aktuelt �
g� utenfor Blindern. Men v�r forsiktig med dette, da for det f�rste
studenter er (og skal v�re) m�lgruppen v�r, og dette kan tiltrekke seg
u�nskede elementer. Enkelte foreninger har pr�vd dette, og droppet det
etter d�rlige erfaringer. Send pressemeldinger til andre aviser/radio/TV etc.
Aktuelle steder kan v�re spalten "Hva skjer?" i Aften Aften, og Dagbladet
og VG har vel tilsvarende. Lag gjerne en liste over stder � sende til,
som kan leveres videre til neste infosjef.


\subsection{Stuntbl�sting}
Med stuntbl�sting menes som oftest � g� p� forelesninger og snakke litt i pausen.
Foreleserne er som oftest positivt innstilt om man sp�r p� forh�nd og ikke bruker
for mye tid. Stuntbl�sting kan ogs� betegne andre aktive opplegg, det er bare �
v�re kreativ. Bi�rnebl�s er ofte interesserte i � v�re med og stuntbl�ste.


\subsection{Andre oppgaver}
Infostyret har ogs� enkelte andre oppgaver i RF, spesielt � produsere
medlemskort og billetter til arrangementer. Medlemskortene m� v�re klare
f�r semesterstart, slik at man f�r forh�ndssolgt ved semsterstartsfesten
og ved rekrutteringsstands tidlig i semesteret (f eks diverse fadderopplegg).

Billettene m� v�re klare tidlig hvis man har tenkt � ha forh�ndssalg
(som ofte kan v�re en god ide ved store arrangementer). Billettene m�
nummereres slik at vi kan se hvor mange som er solgt. Anders Schau Knatten har et skript som genererer numererte billetter klare til kutting.


\subsection{Diverse tips}
Husk at kafeen er et fint sted � bl�ste. Her kan man dele ut l�pesedler
etc. Avtal gjerne med kafesjef at den (eller kafefunksjon�rene) legger ut
l�pesedler p� bordene n�r kafeen �pnes. Dette kan ogs� gj�res p�
kjellerkroer, men det har en tendens til � skape endel rot. Folk tenner
p� l�pesedlene, lager papirfly av dem, heller stearin p� dem etc.

Infosjef b�r ta vare p� alt bl�stmateriell som lages (i elektronisk
form) og gi det videre til neste infosjef, slik at den da kan ta
utgangspunkt i ting som allerede er laget. Det er ikke meningen � bruke
det samme materialet hvert semester, utformingen kan jo v�re helt
anderledes, men det er jo ingen vits i � finne opp hjulet p� nytt.
Man kan jo ogs� f� noen gode ideer p� denne m�ten.

Ved semesterstart i h�stsemesteret kan det v�re lurt � aliere seg med matnat-faddere. Mange av disse er allerede RFere, og b�r ikke v�re noe problem � f� kontakt med. Be fadderne opplyse nye studenter om RF, og gi dem gjerne flyers.

% vim:fileencoding=utf-8
% -*- encoding: utf-8 -*-
% vim:fileencoding=iso-8859-1

\section{Husbjørnen}

Denne teksten er rettet mot de som er med i Husbjørnen. Unntak er dog
den første teksten og den om trykking, som er ment for redaktøren. I
tillegg er teksten om komponering rettet mot personen som er ansvarlig
for design og komponering.

\subsection{Husbjørn-redaktør}
 Husbjørn-redaktøren har som oppgave å produsere, distribuere og stå
til ansvar for innholdet i Husbjørnen - realistenes avis. For å få til
dette er det nesten absolutt nødvendig å ha en redaksjon.

 Veldig mye av arbeidet kan gjøres gjennom dem. Dette gjelder spesielt
innsamling av materiell. Jobben inneholder mange detaljer, og det er
fort gjort å bomme på et eller annet, så vær strukturert.

 Det anbefales sterkt at man holder møter minst en gang i måneden, og
holder en klar saksliste med punkter som skal dekkes, slik at ikke
ting blir glemt. Regn alltid med at ett eller annet blir glemt.

 Sett fristene for innlevering av materiell til minst 14 dager før
bladet går til trykkeriet. Dette hjelper både for personer som er sene
med artiklene sine, og gir tid til personen som designer avisen til å
komponere alt sammen.

 Erfaring viser at hvis tekster blir for lange, så gidder ikke folk å
lese dem. Hold informasjon ned til maksimalt to sider inklusivt
bilder. For A5 formatet, så betyr det en absolutt øvre grense på rundt
1000 ord. Det er da oversiktelig hvis de fleste tekstene er på opptil
rundt 500 ord.

 For å ordne med alt som har med terminalstuer å gjøre, så tar man
kontakt med ansvarlig for Drift: Stein Michael Storleer på
michael@matnat.uio.no. Han kan også gi kontaktinformasjon til de
ansvarlige for de forkjellige terminalstuene. Hvis man skal ha kontakt
med alle stueansvarlige samtidig, kan man eventuelt ta kontakt på
stueansvarlige@matnat.uio.no.

 For å ordne med distribusjon utenom termstuer, så tar man kontakt med
ansvarlig for Teknisk. Dette er for tiden Astri Ottesen med e-post
astri.ottesen@admin.uio.no.

\subsection{Innsamling av materialet}
 Hva som er med er opp til redaktør og redaksjonen. Man kan ikke si at
en måte er bedre enn en annen. Men per dags dato er dette oppsettet
for hva vi gjør:

\subsubsection{Intervju}
 For å lage et godt intervju hjelper det å ha gode spørsmål. For
personlige spørsmål, snakk med mennesker som kjenner ham/henne. For
upersonlige spørsmål snakk med en som kan temaet intervjuobjektet blir
intervjuet for.

 Et profesjonelt intervju gjøres med en diktafon. Hvis dette ikke er
så viktig, anbefaler jeg at man samles foran en Pc og bytter på å
skrive med objektet. Eventuelt kan man til nøds også skrive intervjuet
fortløpende. Dette går dog tregt.

 Få tak i fotografier av intervjuobjektet.

\subsubsection{Pop.vit}
 Forutsagt at man vet hva man ønsker en populærvitenskapelig artikkel
om, så er det bare å ta kontakt med forskjellige professorer på
universitet og spørre. Mange av dem er entusiastiske om tanken om å få
publisert tanker om hva de driver med.

 Få tak i fotografier forfatteren av artikkelen.

\subsubsection{Husbjørnen for $n$ år siden}
 Interessante artikler fra Husbjørnens arkiv. For å få til dette må
man lete i arkivet etter en artikkel man synes er interessant. Dette
må da skrives for hånd inn på en datamaskin siden det ikke finnes
digitale kopier. Redaktør skal ha oversikt over arkivet.

 Eventuelle bilder og ingresser må også ordnes.

\subsubsection{Svada Hall of Fame}
 Sitater fra e-postlisten svada i RF. Dette ordnes typisk gjennom en
fast person som er på denne listen. Hvis man ikke har en, så spør
rundt. Dette er ikke en stor oppgave, så flere er nok villige. Alt man
trenger for å gjennomføre denne spalten er å purre på for sitater,
slik at de kommer inn i tide, og å få bekreftet at man har tillatelser
til å sitere personen på sitatet deres.

\subsubsection{Blindernekskursjonene}
 Vurderinger av forskjellige tilbud rundt omkring på Blindern. Dette
har vært lettest å kombinere som et internarrangement for redaksjonen.
De som blir med får prøve ting gratis, og alle bidrar med idéer til
innhold. Denne ordningen har lykkes ekstremt bra hittil.

\subsubsection{DJ-panelet}
 Diskusjon mellom DJ-ene i RF om forskjellige spørsmål man stiller
dem. Denne spalten ordnes ved å stille et eller annet musikkrelevant
spørsmål direkte på Regis e-postliste: regi@rf.uio.no. Lag et
sammendrag, korriger navnene og send det tilbake til listen for
godkjenning. Typisk gjør folkene på listen veldig mye av jobben selv.

 Legg ved en ingress.

\subsubsection{Formannen har ordet}
 Denne spalten skriver seg selv så lenge man klarer å overtale Formann
til å skrive. Dette er ikke nødvendigvis en enkel oppgave.

 Eventuelle fotografier må også ordnes.

\subsubsection{Dikt}
 Hvis man ikke har et godt dikt på lur, eller kjenner noen som har
det, så kan man igjen ty til Husbjørnens arkiv. Mange av de gamle
avisene har dikt i seg som er både fagrettet og interessante (hvis man
liker den slags).

 Legg ved relevante bilder/fotografier hvis man klarer å skaffe noen.

\subsubsection{Konkurranse}
 Denne står veldig åpen. Så lenge man er villig til å gjennomføre en
bekreftelse av hva enn det er det bes om, så er det få grenser for
spørsmål. Lag gjerne en tekst rundt for å krydre det hele. Husk at det
må være med en premie, en måte å sende inn løsningsforslag til
konkurransen og en klar definisjon av hva det blir spurt om.

 Eventuelle bilder og/eller fotografier må også ordnes.

\subsubsection{Program}
 Denne burde være så oppdatert som mulig. Hold jevnlig kontakt med
Arrangementsjef om dette.

 Lag gjerne noen små tekster til en del relevante arrangementer, slik
at programmet blir mer utfyllende.

\subsubsection{Leder}
 Denne skrives hovedsakelig av Redaktøren, hvis ikke redaktøren sender
oppgaven videre. Den kan brukes som en lufteventil for uttrykke sin
frustrasjon over hvordan ting er.

\subsubsection{Tegneserie}
 Dette er et foretak som krever to elementer: En skribent som lager
plot og dialog, og en tegner som komponerer sammen. (Dette kan være
samme person.) Hvis man ikke har noen til å utføres disse to oppgavene
bra nok, så vil tegneserien falle sammen. Så ikke ha for høye
ambisjoner om innholdet i tegneserien, hvis den ikke har motivasjonen
i orden.

\subsubsection{Ymse}
 Ideer ligger og virrer rundt omkring overalt. For å få realisert dem,
så trenger man bare å få samlet en konkret vinkling, og å spørre rundt
om noen kunne tenke seg å skrive denne. Så lenge oppgaven er konkret
nok og den virker interessant, så er det flere i foreningen som kan
fungere som ypperlige frilansskribenter. Dessuten skal man også alltid
vurdere om en tekst kan konverteres til en spalte slik at man får mer
ut av den.

 Husk å legge ved ekstra informasjon som fotografier og/eller bilder
hvis det passer.

\subsubsection{Forsidetegning}
 Dette er typisk en oppgave man venter med til rundt deadline før man
starter på. Det er ikke før da man har nok stoff til å si hva du har
lyst til å ha på forsiden. Enkleste løsningen er å be noen om å tegne
noe som samsvarer med en av tekstene på innsiden. En bjørn er dog en
klassiker som sjeldent slår feil.

\subsection{De store OBS-ene}
 Noen av de største feilene som kan oppstå i en avis, er at man bruker
tekster som enten er feil, eller som man ikke har tillatelse til å
bruke. Førstnevnte er ikke et veldig stort problem, siden Husbjørnen
sjeldent publiserer fakta. Sistnevnte derimot skal man være litt
forsiktig med.

 Hvis man skal ha et fotografi av en person, husk å be om tillatelse
fra personen før man bruker det.

 Hvis man skal sitere noe, husk å spørre om lov, og å markere
hvem/hvor det kommer fra i teksten på en eller annen måte.

 Å vite om et arrangement er offentlig eller internt kan ofte være
vanskelig å bedømme. Spør alltid arrangøren om lov, hvis man skal
publisere oppføringen av et arrangement.

 Mange bilder på nett kan det være farlig å kopiere av. Den beste
måten å unngå det på er å få en tegner til å tegne en kopi. Da er man
på en tryggere side av loven. Unntaket er selvfølgelig hvis man har
bedt rettighetshaverne om lov først.

 Skal man referere til et navn eller en e-postadresse, be om lov.

 Det finnes en e-postliste som heter korrrektur@rf.uio.no. Den
e-postlisten er dedikert til å korrekturlese alle tekster som sendes i
deres retning. Jeg advarer mot å bruke listen til personlige formål,
for å ikke irritere de som står på listen for mye. Hvis det finnes en
eneste mistanke om at en tekst på en eller annen måte har bruk for
korrektur, send det til denne listen. Klipp og lim teksten som skal
korrekturleses rett inn i en e-post og still din forespørsel på toppen
av teksten. De er veldig lite glad i vedlegg.

 Når en tekst er ferdig, skal den enten til redaktør eller til
designer. Dette kommer an på om redaktør ønsker å overse ting før det
publiseres eller ikke. Igjen skal teksten være i formen av en e-post
uten vedlegg.

 Bilder skal ikke sendes som vedlegg. De skal enten overføres direkte
til enten redaktør eller designer, eller gjennom en link rettet mot
hvor filen er lagret på nettet.

\subsection{Komponering}
 For å komponere avisen brukes programmet InDesign. Denne finnes i
flere versjoner, og det kan være farlig om to versjoner er
inkompatible, og flere enn en skal ta en titt på samme utgave.
Programmet er relativt intuitivt, og støttes lett opp mot Photoshop.

 Hvordan avisen skal se ut er mer eller mindre opp til designeren. Det
er dog viktig å ha i bakhodet at dette er avisen til
Realistforeningen, og at man ikke skal fravike for mye fra tema og
logoer. Typisk skal det alltid være en del bjørner med i bladet.

\subsection{Trykking}
 Trykking skjer typisk gjennom CopyCat. Ved å sende en e-post til
nydalen@copycat.no så er mesteparten av jobben gjort. Introduser deg
selv, og be om et prisoverslag over trykkejobben du ønsker å få
utført. De fører ikke prislister, siden kalkuleringen av trykking er
en noe komplisert affære. Dessuten kan man diskutere prisen noe, hvis
man ber om en veldig stor jobb. Derimot kan man be om flere
prisoverslag om gangen. Pass på at de regner prisen uten moms. De
trenger å vite:

\begin{itemize}
 \item Format: Typisk A4 eller A5
 \item Vekt: Standard ark er på 80g eller 100g. Mer eksklusive får man ved
100g glanset og 200g glanset.
 \item Farge: Hvis man skal ha noen av arkene i en annen farge, eller man
ønsker fargesider, så må det spesifiseres.
 \item Antall: Hvor mange eksemplarer som skal trykkes.
\end{itemize}

 Per dags dato kjører vi A5, 100g glanset utside, 80g innmat og 400 eksemplarer.

 Hvis man er usikker på om hvordan ting vil se ut i et format man ikke
har sett på så kan man be om å få se på prøveeksemplarer. Dette koster
ikke noe.

 Når man skal sende ting til trykkeriet, så er det bare å sende en
e-post med avisen i form av en PDF-fil vedlagt. Eventuelt kan man
sende en link til et sted hvor man har publisert filen på nett.

 Hvis man ønsker å vite hvor lang tid det tar å trykke, så er det lite
som er i veien for å spørre. Typisk tar det en dag.

\subsection{Distribusjon}
 Når avisene er ferdig trykket, så kan de plukkes opp hos Copycat.
Dette ligger i Nydalen, på venstre side av BI bygget sett fra Nydalen
t-banestasjon. Det tar ca. to minutter å gå fra stasjonen.

 Ordningen per i dag er at det alltid er to personer som går sammen.
Dette gjøres av to grunner. Den ene grunnen er at det blir mer sosialt
å gjøre jobben. Den andre grunnen er at hvis bare en person vet hva
som skal gjøres, vil den andre også lære, og kan spre sin kunnskap
videre.

 Per i dag så har Husbjørnen et gitt sett med steder hvor den kan
plasseres. Man kan ignorere tillatelsene som er satt, men man
risikerer at bladet kan bli fjernet. Flere steder er ikke helt i orden
heller, siden vi trenger avisstand for at ting skal være helt korrekt.
På sikt er det målet. Inntil videre kan man prøve snike avisen inn ved
siden av andre magasiner hvor det er plass.

 Budrunden tar drøye to timer, og etter stengetid så krever det også
ofte tilgangskort. Pass på at dette er i orden før man begynner å
distribuere. Husk også å samle opp alle gamle blader. Disse skal
redaktøren ha.

 Vi har spredt ting ut på stands og terminalstuer rundt omkring på
Blindern. De med en «*» ved mangler per dags dato stand.

\begin{itemize}
 \item Biologi. Ved inngangen nærmest PO
 \item Biologi. I utskriftshyllen i terminalstuen plassert lengst unna PO
 \item PO. I utskriftshyllen ved printeren ved inngangen.
 \item Sophus Lie. Ved inngangen
 \item Nils Henrik Abels hus. I utskriftshyllene i terminalstuen i kjelleren
 \item * Fysikk. Ved pendelen
 \item Fysikk. I utskriftshyllen i terminalstuen i fjerde etasje.
 \item * Kjemi. Utenfor terminalstuen. (Foreløpig i terminalstuen på
rotebordet ved inngangen)
 \item * IFI. Ved inngangen av Terminalstuen i første etasje. (foreløpig
på rotebordet)
 \item * IFI. Ved inngangen til kantinen i andre etasje.
\end{itemize}

 På IFI har vi fått klar beskjed om at «det er bedre å be om
tilgivelse enn tillatelse». Dvs. at vi har ganske frie hender der.



\chapter{Desisjonsutvalget}
% vim:fileencoding=utf-8
% -*- encoding: utf-8 -*-
% vim:fileencoding=latin1
% -*- encoding: iso-8859-1 -*-
Fra foreningens lover, §14:
\begin{description}
\item[d] Desisjonsutvalget har den endelige avgjørelse i tvilsspørsmål om
tolkning av lovene. Utvalget kan også fatte vedtak og gi regler i situasjoner
hvor lovene måtte vise seg å være utilstrekkelige. Ethvert medlem av
Realistforeningen har rett til å innanke for Desisjonsutvalget vedtak hvor det
kan være tvil om lovligheten.
\item[e] Ved mistanke om misligheter kan Desisjonsutvalget suspendere
medlemmer av styrer, utvalg og komiteer. Suspensjonen kan omfatte et organ i
sin helhet, selv om det ikke foreligger konkret mistanke mot hvert enkelt
medlem. I tilfelle suspensjon er foretatt skal Desisjonsutvalget straks sørge
for at det blir innkalt til ekstraordinær generalforsamling der
mistillitsforslag behandles og nyvalg eventuelt avholdes.
\item[f] Desisjonsutvalget har ansvar for at Realistforeningens arkiv til 
enhver tid er i orden.
\item[g] Desisjonsutvalget har møte- og talerett i alle foreningens organer.
\end{description}

\section{Lovtolking}
RF er, som kjent, et monarki, og vi praktiserer liten grad av sedvanerett. Derfor har vi
heller ingen fast tolkning av lovene utover hva noen klarer å huske. Tolking av lovene
må derfor skje hver gang noen stiller spørsmål ved noe. Desisjonsutvalget har ved enkelte
anledninger funnet meget romslige tolkninger, men de har likevel passet med lovteksten.

Når det er nødvendig, bør lovene tolkes strengt. Desisjonsutvalgets oppgave er dog å passe på
at foreningen fungerer godt, ikke stikke kjepper i hjulene for styret. Man kan for eksempel
tenke seg at økonomiske rutiner er viktigere enn andre ting.

\subsection{Lovendringer}
Desisjonsutvalget skal uttale seg om lovendringsforslag som er kommet opp på Generalforsamling.
Det er her viktig at Desisjonsutvalget gjør et ordentlig forarbeide, for Desisjonsutvalgets
anbefalinger blir nesten alltid tatt til følge.

Lovendringsforslag skal komme inn senest en uke før Generalforsamling. Etter dette har Desisjonsutvalget
adgang til å komme med forslag inntil 48 timer før. Denne adgangen bør kun brukes til å forbedre
innkomne lovendringsforslag eller til å lage kompromissforslag når det er kommet inn motstridende
forslag. Desisjonsutvalgets medlemmer bør ellers overholde den vanlige fristen for lovendringer.

\section{Suspensjon}
Denne adgangen er gitt Desisjonsutvalget fordi der finnes en kuppekultur blant 
studentforeningene. Desisjonsutvalget kan, hvis for eksempel SVFF skulle kuppe en Generalforsamling,
kaste styret og kalle inn til ny Generalforsamling. Ofte vil det heller være mulig å tolke
medlemsparagrafen på Generalforsamling. Det kan også tenkes andre tilfeller hvor denne adgangen
kan benyttes. Et eksempel er underslag.

\section{Arkivet}
Arkivet ligger i Jomfruburet. Dette er et rom plassert i 5.~etage i fysikkbygningen. For å komme
dit, trenger man en nøkkel. Denne nøkkelen passer til store fysiske lesesal, den østlige døren i 
nordveggen der og til selve Jomfruburet, som er en etage opp trappen fra denne dør.

\section{Møterett}
DU kan stille på alle møter. Av den grunn er det viktig at alle innkallelser kommer i tide. Desisjon
møter av og til når det kan være tvil om lovlighet eller lovspørsmål skal tas opp.

\section{E-post}
Listen desisjon@rf.uio.no skal stå som metamedlem på alle styrelister: as@rf.uio.no, hs@rf.uio.no, ks@rf.uio.no,
blaestgruppa@rf.uio.no, fagstyret@rf.uio.no, gildet@rf.uio.no, okonomi@rf.uio.no og revisjon@rf.uio.no. Desisjonsutvalgets
medlemmer bør unngå å blande seg opp i daglig drift, men det er greit å forklare hvordan ting var før og andre
faktaopplysninger. Disse opplysningene bør være av nøytral karakter.



\chapter{Andre verv}
Verv som ikke er direkte knyttet til bestemte styrer eller utvalg.
\section{Internansvarlig}

\subsection{Mandat}
Internansvarlig har som oppgave � legge til rette for et godt sosialt milj�
i foreningen, og f� folk til � trives. Vervet inneb�rer � gj�re seg kjent
med alle aktive interne medlemmer, med spesiell vekt p� nye medlemmer,
og s�rge for at de trives gjennom hele semesteret.

Dette inneb�rer ikke bare at de trives sosialt, men ogs� at de mestrer sine
egne verv. Internansvarlig skal derfor ogs� holde et �ye �pent for om noen
trenger mer oppl�ring eller hjelp til arbeidsoppgavene, og i s� fall ta
initiativ til at de f�r dette.

Internansvarlig har taushetsplikt dersom interne som tar kontakt om et
hvilket som helst tema �nsker dette.

Internansvarlig har ansvar for � oppdatere denne teksten med nye erfaringer etter fullf�rt verv.


\subsection{Typiske oppgaver}
En viktig praktisk oppgave vil v�re � holde oppdaterte lister over interne til
enhver tid, med fullt navn, telefonnummer, epostadresse og postadresse.

Det er viktig at internansvarlig holder god kontakt med kjellerstyret,
som har de fleste funksjon�rene under seg.

Ellers er internansvarlig er et ganske fritt verv, hvor oppgavene m� tilpasses
behovet i foreningen og enkeltpersonene. Det er viktig at der andre ser
p� hvordan vervet skj�ttes, ser internansvarlig p� hvordan personen bak
vervet trives. Typiske oppgaver vil v�re � jobbe for � opprettholde
Mandagstradisjonen, s�rge for at internfestene ikke blir glemt,
og v�re med og arrangere andre internarrangementer som hyttetur og
Danmarkstur. Det er spesielt viktig � invitere nye interne med p�
internarrangementene, slik at de trives og blir i foreningen.

SiO L�ringsmilj� holder p� med � utvikle et hefte om utbrenthet.
Det vil v�re naturlig at Internansvarlig henter r�d derfra og
jobber for � forebygge dette i foreningen.

\subsection{Tips og triks}
Internansvarlig er forel�pig et verv som ikke har v�rt bemannet
p� en stund, og det er derfor opp til den som tar vervet � pr�ve
seg frem og gj�re seg erfaringer.

Det er i hvert fall viktig at internansvarlig er tilstede ved
kranf�rerkurs og andre rekrutteringstiltak, for slik � bli kjent
med nye interne. Det kan ogs� v�re en fordel om internansvarlig
hilser p� alle nye funksjon�rer i baren, kafeen og bl�st og
sp�r hvordan de trives.

\subsubsection{Om informasjon}
E-post er en viktig informasjonskanal, s�rg for � bruke denne
aktivt til � informere om alt som skjer i god tid. Men e-post er
et tveegget sverd, da mange ikke leser denne i det hele tatt eller s�rlig
ofte. Det er derfor viktig � bruke andre informasjonskanaler som
jungeltelegrafen, SMS (det jobbes for tiden med � f� til en avtale om bruk
av SMS-server til masseutsendelser) og telefon. Og ikke minst, snakk
med s� mange som mulig personlig. Dette er spesielt viktig ovenfor
nye interne, som ikke kjenner milj�et og tradisjonene s� godt.
\section{Donaldsjef}
Donaldsjef har ansvar for vedlikehold og oppdatering av Realistforeningens store Donaldsamling.


\chapter{Generalforsamling}
% vim:ft=tex:fileencoding=utf-8:tabstop=4:shiftwidth=4
% -*- encoding: utf-8 -*-

\section{Forberedelser}
Det er ikke mye som må gjøres foran en ordinær Generalforsamling, men det er
viktig at ting blir gjort.

\subsection{Arrangement}
Tidligere har man kombinert Generalforsamling med andre aktiviteter. Dette kan 
være en idé for å trekke folk. Generalforsamlingen bør blæstes godt for Interne,
men alle betalende medlemmer har adgang.

Det har vært vanlig å sette frem noen kasser Øl til de som dukker opp, og etter
Generalforsamlingen er det gjerne pizza på huset. Det er ikke sikkert man blir
settedyktig, så vent med serveringen til man er igang.

\subsubsection{Oppgavefordeling}
Det er lurt å ha kandidater til formalitetene klart på forhånd. Formannen er
tradisjonelt ordstyrer, og bør ikke ha andre oppgaver under møtet. Sekretær
kan med fordel skrive protokoll. Personer som underskriver protokoll bør være
mulig å få tak i, og bør finnes på forhånd. Dette bør helst være andre personer
enn Formann, Sekretær eller DU.

Videre: Pizza må bestilles, og øl kan kjellermester skaffe.

På forhånd er det lurt å ta kontakt med Regi, slik at det blir vist Donald-filmer
i pausen(e). Auditorium 1 i VB er vanligvis lokalet; dette må holdes av.

\subsection{Innkalling}
Det må kalles inn til Generalforsamling minst ti virkedager før. Feriedager teller
ikke som virkedag. \emph{I tillegg er det verdt å merke seg at virkedager telles på sekundet:
Er generalforsamlingen berammet til onsdag klokken 18.04, skal innkallelsen sendes ut
(hvis det ikke er ferier i veien) senest klokken 18.04 fjorten dager tidligere.} I skrivende
stund er lovene slik at man må være medlem når inkallelsen kommer for å få stemmerett, og derfor
er det lurt å kalle inn etter siste arrangement før fristen.

Det er også lurt å legge forslag til dagsorden sammen med innkallelsen. Forslag til dagsorden
må senest sendes ut fem virkedager før møtet, men det er ingen grunn til ikke å ta det med
allerede ved innkallingen. Det er viktig å merke seg at Generalforsamlingen bare kan fatte vedtak
når saken er foreslått tatt opp fem virkedager før møtet.

Innkallelsen må offentliggjøres. En epost til informasjon@rf.uio.no holder til det, men det 
bør også henges opp papirkopier på RFs tavler. Innkallelsen bør også inneholde et notat
om listen diskusjon@rf.uio.no og hvordan man melder seg på denne – hvis det varsles på 
innkallelsen, kan denne brukes til å foreslå saker for Generalforsamlingens overveielse,
inkludert lovendringer. Innkallelsen bør isåfall inneholde korrekt prosedyre for å foreslå
saker.

\subsection{Økonomi}
Hovedstyrets primæroppgave er å tenke langsiktig, og i så måte er budsjettet viktig.
Hovedstyret skal innen en uke før Generalforsamlingen komme med et forslag til budsjettrammer
som Generalforsamlingen skal vedta. Det er viktig at dette forslaget er lettleselig og godt
kommentert, for andre personer skal kunne fremme endringsforslag innen 48 timer før møtet.
Det er viktig at det finnes en reell mulighet for dette, for fremtidige planer avhenger
gjerne av budsjettet.

I tillegg må Økonomiutvalget kunne gjøre rede for RFs økonomi, og Revisjonsutvalget
må skrive sin beretning.

\subsection{Semesterberetning}
Semesterberetningen kommer gjerne skriftlig med sakspapirene. Hver gruppering skriver sin.
Få med navn på interne og hva som er gjort dette semestret, gjerne med en liten evaluering.

\section{Lovendringsforslag og andre saker}
Hvis Generalforsamlingen skal fatte vedtak i en sak, må saken foreslås tatt opp minst
fem virkedager før. Forslagsstiller må selv ta ansvar for å offentliggjøre forslaget, f.eks.
ved å sende forslaget til diskusjon@rf.uio.no hvis dette er foreslått på innkallelsen.
Hovedstyret (hs@rf.uio.no) skal ha kopi av alle forslag.

Hvis forslaget er et lovendringsforslag, er det strengere krav enn andre forslag. Et lovendringsforslag
må leveres inn med alle formuleringer: Man kan ikke bare foreslå en intensjon, og overlate til
Generalforsamlingen eller andre å formulere den. Generalforsamlingen kan ikke endre på ordlyden
i lovendringsforslag, bare vedta eller avvise forslaget.

Saker det ikke skal fattes vedtak i, kan foreslås når som helst. Men merk at disse begrensningene
er til for at ikke Generalforsamlingen skal ta for mye tid – en ordinær Generalforsamling tar 
allerede to–tre timer, og for å bli vedtaksdyktig trenger man gjerne en del mennesker som
ikke er interessert i detaljer omkring lovverk.

\subsection{Desisjonsutvalget}
Desisjonsutvalget har utvidede rettigheter til å foreslå endringsforslag. Hvis desisjonsutvalget
liker forslaget, er det vanlig praksis at det retter opp «formaliafeil» eller foreslår en annen ordlyd.
Desisjonsutvalgets muligheter gjør det også mulig å rette opp konflikter i innsendte forslag.

\section{Avstemninger}
Det er forskjellige former for avstemninger. RF bruker alminnelig flertall, absolutt flertall
og 2/3-flertall.

\subsection{Alminnelig flertall}
Alminnelig flertall har det alternativ som får flest stemmer. Alminnelig flertall er den vanligste
formen på generalforsamlingen, og er også formen som brukes ved akklamasjon.

For alminnelig flertall, holder det med én stemme så lenge ingen stemmer noe annet. Derfor kan mange
avstemninger gjennomføres ved først å forsikre seg om at ingen stemmer mot, for hvis ordstyrer stemmer
for, er denne stemmen nok.

Hvis det er like mange kandidater som verv (eller bare ett alternativ å stemme over), kan man gjennomføre
avstemningen ved akklamasjon – å klappe betyr å stemme for.

Hvis det er flere kandidater, stemmer de stemmeberettigede på like mange kandidater som det er verv.
De som har fått flest stemmer, blir da valgt.

\subsection{Absolutt flertall}
Hvis det skal velges én person til et verv, kan tre personer kreve absolutt flertall. Absolutt flertall
betyr at minst halvparten av forsamlingen må stemme for alternativet.

Skriftlig valg hvor det er krevd absolutt flertall, kan med fordel avgjøres etter preferansevalgmetoden.
I teorien betyr dette at man har flere valgomganger, og fjerner den minst populære kandidaten i hver omgang
inntil noen har absolutt flertall. Preferansevalg innebærer at kandidatene rangeres på stemmeseddelen,
og at tellekorpset teller stemmesedlene flere ganger og gir stemme til den første kandidaten som fortsatt er med
for hver seddel.

En annen metode for absolutt flertall, er «majoritetsvalg», der de to mest populære kandidatene går videre til
en ny runde hvis ingen får absolutt flertall ved første forsøk. Denne metoden brukes f.eks. i det franske presidentsvalget.

Hvis flere skal velges samtidig, kan absolutt flertall ikke benyttes.

\subsection{2/3-flertall}
Lovendringer vedtas med 2/3-flertall. Dette krever at 2/3 av forsamlingen aktivt må stemme for forslaget.
Dersom man har flere motstridende forslag, bør man derfor først finne den mest populære ved å sette to og to
forslag opp mot hverandre, og ta avstemning for vinneren.

\section{Møtets gang}
Generalforsamlingen behandler mange saker. Noen er vedtakssaker, mens andre er orienteringssaker.
I vedtakssaker \emph{må} referatet gjengi vedtaket. Dette er de typiske saker, ikke alle står i lovene.

Generalforsamlingen settes ikke før man er settedyktig, og da bæres HKH Ursus Minor inn (alle reiser seg), og
han får en øl.

\subsection{Godkjennelse av innkallelsen}
Innkallelsen kan bare godkjennes hvis Desisjonsutvalget konkluderer med at den var gyldig.
Dersom så er tilfelle, kan generalforsamlingen vedta å godkjenne. Ordstyrer bør her spørre
om noen tegner seg imot innkallelsens godkjennelse. Hvis ingen tegner seg mot, kan ordstyrer
– hvis han har stemmerett og stemmer for godkjennelse – trygt anta at dette har flertall.
Hvis noen stemmer mot, må flere stemme for.

\subsection{Godkjennelse av dagsorden}
Generalforsamlingen kan endre dagsorden frem til den er vedtatt, men det er ikke anledning til å
fremme nye vedtakssaker. Dagsorden vedtas typisk på samme måte som innkallelsen.

\subsection{Valg av ordstyrer og referent}
Ordstyrer og referent er gjerne hyret inn på forhånd, og velges som regel ved akklamasjon.

\subsection{Valg av to personer til å underskrive protokollen}
Disse er helst spurt på forhånd. Ellers bør ordstyrer foreslå passende personer fra forsamlingen.

\subsection{Valg av to personer til å plukke papirfly}
Hvis/når HKH Ursus Minor treffes av et papirfly, foreslås denne personen til papirflyplukking.

\subsection{Regnskaper}
ØU-leder får gjerne innlede denne. Regnskapet presenteres (det er fint om en kopi er delt ut) på overhead e.l.,
og forsamlingen bør få en forståelse for foreningens økonomi.

Revisjonsutvalget vil legge frem regnskaper som er revidert med tilhørende revisjonsberetning. Etter
RUs innstilling, kan generalforsamlingen med simpelt flertall godkjenne eller ikke godkjenne regnskapet.
Godkjennelse av regnskapet betyr at de ansvarlige innvilges ansvarsfrihet for regnskapet.

\subsection{Budsjettrammer}
Alle forslag til vedtak, må være offentliggjort innen generalforsamlingen. Et av dem velges med simpelt 
flertall.

\subsection{Lovendringsforslag}
Desisjonsutvalget vil gjerne kunne informere om forslagene. Et forslag må vedtas med 2/3-flertall.
Det vil ofte være megen debatt ved lovendringsforslag. Ordstyrer må her sørge for at dette ikke går over
styr. Det kan være lurt å ha hjelp til å opprettholde en taleliste. Pass på at denne følges. Hvis 
argumenter gjentas, kan taleren stanses (klubbes).

\subsection{Semesterberetninger}
Dette er en orienteringssak/diskusjonssak. Semesterberetningene bør deles ut på forhånd, og medlemmene
kan stille spørsmål til de interne. Det foretas ingen avstemning under denne saken.

\subsection{Kontigenter}
Forslag til kontigent fremmes på Generalforsamlingen. Et av forslagene velges med simpelt flertall.
Hovedstyret kan med fordel komme med et forslag i sakspapirene. Merk at det ikke finnes noe underforstått
forslag.

\subsection{Valg av tillitsvalgte}
De som stiller til valg får presentere seg, og deretter stemmes det. Det er vanlig å kreve skriftlig valg
dersom det er flere kandidater enn verv. Isåfall bør det oppnevnes et tellekorps. Under valg er det vanlig
å ta hensyn til at noen kanskje stiller til flere verv, og rekkefølgen på valgene kan justeres etter dette.

\subsection{Ymse}
Ymse kan brukes til å fortelle om viktige saker, og å spørre om detaljer. Generalforsamlingen kan ikke
fatte noe vedtak.



\appendix

\chapter{Oppdatering av denne boken}
Det er ønskelig at denne boken hvert semester oppdateres i henhold til
de erfaringer man har hatt. Boken er skrevet med \LaTeX, og det er ikke
sikkert alle har kjennskap til dette systemet. Det kan virke komplisert,
men for å holde dette dokumentet oppdatert trenger man svært lite
kunnskaper.

\LaTeX{} kan man finne ut mer om på \url{http://www.latex-project.org}. Termstuene
har også bøker om \LaTeX{} til utlån.

Hviteboken holdes under kontroll med revisjonskontrollsystemet Subversion.
Dette programmets hjemmeside finnes på \url{http://subversion.tigris.org/}, og
der kan også alle programmer man trenger lastes ned.

Hviteboken ligger på \url{https://sub.ifi.uio.no/repos/users/rf-hvitebok} under katalogen trunk.
Man må ligge i filgruppen til RF for å få tilgang. Tilgang kan man få ved å 
kontakte Hovedstyret. Hovedstyret gir tilgang ved å logge seg på 
\url{https://wwws.ifi.uio.no/system/svn/} med brukernavn rf.

\section{Skaffe tilgang til filene}
For å jobbe med filene, trenger man en lokal arbeidskopi. I Linux\footnote{Eller et annet UNIX-lignende system med subversion og ssh} kan man
skaffe dette med kommandoen\\
\verb@svn co --username <bruker> https://sub.ifi.uio.no/repos/users/rf-hvitebok@\\
i xterm/terminal. Erstatt «<bruker>» med egnet UiO-brukernavn. Passord er da uio-passord.

\section{Oppdatere eksisterende filer}
Hoveddokumentet heter «hvitebok.tex». Under katalogen «kildekode» ligger
resten av filene, stort sett en fil for hvert verv pluss noen til.
Under katalogen «kildekode/arr» ligger arrangementsbeskrivelser og
erfaringsprotokoller ifm arrangementer. Alle disse filene er tekstfiler
som kan åpnes i et hvilket som helst tekstredigeringsprogram.
(Filen «Makefile» inneholder regler som benyttes av kommandoen «make»
for å produsere dokumenter av forskjellig format ut fra kildekoden).

\section{Legge til filer}
Legg filen i en av katlogene beskrevet over. Filen må være ren tekst,
og kalles «filnavn.tex». Bruk samme koder for overskrifter og lister
som du finner i de eksisterende filene. Legg til filnavnet i «Makefile»,
fremgangsmåten er beskrivet i den fila. Legg til filnavnet på riktig
sted i «hvitebok.tex», på samme måte som alle de andre filene er skrevet.

Nye filer må legges til i subversion med kommandoen \verb@svn add <filnavn>@.
Først da vil Subversion bry seg om filen. Vær vennlig og ikke legg inn filer
som genereres av programmer. For eksempel hvitebok.pdf skal ikke inn i Subversion.

\section{Oppdatere databasen}
Når du jobber med din lokale kopi, vil ikke endringene vises for andre før man
sender inn endringer. Dette gjøres ved å stå i rotkatalogen for prosjektet (katalogen
som inneholder hvitebok.tex) og gi kommandoen \verb@svn ci@. Denne vil gi deg en teksteditor
hvor du kan legge inn en kommentar, som kan gjøre det enklere å finne ut hvilken revisjon
som har hva og hva som er endret siden sist. Dette gir en god changelog\footnote{svn log}.

\section{I Windows}
På termstuene er tortoisesvn installert. Med dette programmet, kan du oppdatere på denne måten:
1) Lag en ny mappe. 2) Åpne denne i explorer. 3) Høyreklikk i vinduet, og velg <<svn checkout>>.
4) Skriv inn \url{https://sub.ifi.uio.no/repos/users/rf-hvitebok/} i URL-feltet. 
5) Benytt UiO-brukernavnet/passordet ditt i feltet.

Filer kan legges til ved å høyreklikke dem, og velge <<add>> fra <<tortoisesvn>>-menyen. Når man er
ferdig med redigeringen, man man høyreklikke mappen og velge <<svn commit>>, fylle ut info, og så
velge <<svn update>>.

\section{Lage pdf- og ps-filer}
Skriv «make» eller «make LaTeX=ltx» på IfI.
Evt «make logclean» for å slette logger.
Andre kommandoer du kan bruke er «make pdf» og «make ps».

% vim:fileencoding=utf-8
% -*- encoding: utf-8 -*-
% -*- encoding: latin1 -*-
% vim:fileencoding=latin1
\chapter{Alkoholreklame \textendash{} Hva er lovlig/ulovlig reklame (i praksis)}
\section{Annonser for skjenkesteder}

Et skjenkested har adgang til å annonsere med navn, beliggenhet, bevillingsrettigheter og åpningstider.

Alkoholprodusenters firmanavn/firmamerke kan ikke forekomme i annonsen.

Navn på eller avbildninger av alkoholholdige drikker kan ikke forekomme i annonsen.

Prisopplysninger må ikke forekomme i annonsen. Som prisopplysninger anses også begreper som «Happy Hour»,«Byens billigste» eller liknende betegnelser som informerer om prisene/prisnivået på alkoholholdige drikker.

\section{Plakater og skilt i tilknytning skjenkested}

Utendørs og innendørs plakater, skilt, lysreklame og vindusreklame kan ikke være merket med alkholprodusenters firmanavn/firmamerke, navn på alkoholprodukter eller avbildninger av alkoholholdige drikker

\section{Prisopplysninger på skjenkesteder}

Som en konsekvens av kundenes krav på å få vite hva de enkelte varene koster, er det tillatt å opplyse om prisene på alkoholholdige drikker. Reklameforbudet begrenser imidlertid på hvilken måte og i hvilken form disse prisene kommer til uttrykk.

Utenfor skjenkestedet kan det gis generell prisinformasjon uten å fremheve spesielle produkter, f. eks. øl fra kr 30\textendash60,\textendash

På skjenkestedet kan priser på alkoholholdig drikk bare gis på lik linje med prisinformasjon for andre produkter.

Kravet om at priser på alkoholholdig drikk bare kan gis på lik linje med annen prisinformasjon innebærer at det ikke er tillatt å benytte iøynefallende eller fremtredende

plakater eller effekter for slike prisopplysninger. Det er heller ikke tillatt å bruke suggererende eller salgsfremmende metoder.

Prisopplysninger på alkoholholdige drikker kan bare komme til uttrykk på en nøytral og ren informativ måte. I motsetning til hva som gjelder for andre varer, er det ikke tillatt å opplyse om tilbudspriser, prisavslag eller reduserte priser på alkoholholdige drikker.

Prislister inne i skjenkelokalet kan inneholde alkoholprodusentens firmanavn, for eksempel «1/2 liter Ringenes pils kr 45,\textendash», «1/2 liter Hansa pils kr 45,\textendash».

\section{Gjenstander på skjenkesteder}

I utgangspunktet vil merking av gjenstander med alkoholprodusenters firmanavn/firmamerke samt navn på/avbildninger av alkoholholdige drikker anses som ulovlig reklame for alkoholholdig drikk.

Av hensyn til den praktiske gjennomføringen av alkoholserveringen, er det imidlertid gjort særlige unntak for vanlig serveringsutstyr. Se avsnittet «Unntak fra hovedregelen».

\section{Plassering av alkoholholdig drikk i skjenkelokalet}

Plassering av alkoholholdig drikk er tillatt så lenge denne kan begrunnes ut i fra hensynet til den praktiske gjennomføringen av serveringen.

Dersom skjenkelokalet er lite eller har begrenset lagerkapasitet, er det en videre adgang til å plassere alkoholholdig drikk i selve skjenkelokalet, enn for skjenkesteder med store lagringsmuligheter.

Hvorvidt plasseringen er lovlig må avgjøres i hvert enkelt tilfelle. Ved vurderingen må hensikten med plasseringen og individuelle forhold på stedet tas i betraktning.

I tilfelle hvor alkoholholdige drikker er plassert på en slik måte at den går over til å være spesiell utstilling, anses dette som ulovlig reklame. Generelt kan spesiell utstilling beskrives som en ekstraordinær fremheving av alkoholprodukter. Selv om plasseringen forklares med at den er til pynt, hygge eller dekorasjon, er den likevel ulovlig.

Reglene for annonser, plakater, skilt, plassering og prisopplysninger på salgssteder er i hovedsak identiske med de regler som gjelder for skjenkesteder.

\section{Prisopplysninger}

Opplysninger om priser på alkoholholdige drikker kan bare gis inne i salgslokalet.

Spesielle regler gjelder for salg via internett. Se avsnittet «Unntak fra lovens hovedregel».

Opplysninger om priser på alkoholholdige drikker kan bare gis på lik linje med prisopplysninger for andre varer.

Prislister inne i salgslokalet kan inneholde alkoholprodusenters firmanavn for eksempel: «0,33 l Ringnes pils kr 15,\textendash»,«0,33 l Hansa pils kr 15,\textendash».

Prislistene skal være utformet og plassert på samme måte som prislister for salgsstedets øvrige varer. Det er ikke tillatt å annonsere eller på annen måte tilkjennegi reduserte priser eller tilbudspriser på alkoholholdige drikker, verken i annonser eller via plakater i eller utenfor salgsstedet.

Prislister for alkoholholdig drikk kan ikke inneholde alkoholprodusenters firmamerke.

\section{Plakater og skilt plassert utenfor salgslokalet}

Skilt/plakater utenfor salgsstedet kan opplyse om at det selges øl på stedet. Opplysningen om at det selges øl skal komme til uttrykk på en nøytral og informativ måte. Det er ikke tillatt å opplyse om hvilke ølmerker som selges på stedet.

Skilt/plakater plassert utenfor salgslokalet kan ikke inneholde prisopplysninger eller alkoholprodusenters firmanavn/firmamerke.

Skilt/plakater utenfor salgssted kan ikke inneholde navn på bestemte alkoholholdige drikker eller avbildninger av slike.

\section{Plakater og skilt plassert inne i salgslokalet}

Skilt/plakater inne i salgslokalet kan inneholde informative opplysninger om hvilke ølmerker (alkoholprodusentens navn) som selges på stedet og prisene på disse uten at noen produkter fremheves i forhold til andre.

Andre typer innendørs plakater, skilt, lysreklame og lignende, kan ikke merkes med alkoholprodusenters firmanavn/firmamerke, navn på alkoholprodukter eller avbildninger av alkoholholdige drikker.

Alkoholreklame kan heller ikke forekomme på kjøleskap eller liknende gjenstander.

\section{Plassering av alkoholholdig drikk i salgslokalet}

Plassering av alkoholholdig drikk er tillatt så lenge denne kan begrunnes ut i fra hensynet til den praktiske gjennomføringen av salgsvirksomheten.

Hvorvidt plasseringen er lovlig må avgjøres i hvert enkelt tilfelle. Ved vurderingen må hensikten med plasseringen og individuelle forhold på stedet tas i betraktning.

I tilfelle hvor alkoholholdige drikker er plassert på en slik måte at den går over til å være spesiell utstilling, anses dette som ulovlig reklame. Generelt kan spesiell utstilling beskrives som en ekstraordinær fremheving av alkoholprodukter. Selv om plasseringen forklares med at den er til pynt, hygge eller dekorasjon, er den likevel ulovlig.

Plassering som tydelig tar sikte på å bli sett utenfra er ikke tillatt. 


\chapter{Historie}
Hadde vært fint med litt historie her.

\chapter{Realistforeningens lover}
RFs lover, oppdatert etter hvert semesters Generalforsamlings 
endringer, skal ligge ved Hviteboken. Dersom noe av innholdet
i Hviteboken ikke samsvarer med lovene, er det selvfølgelig
lovene som gjelder. Dersom gjeldene praksis er i strid med
lovene, må det enten fremmes lovendringsforslag, eller praksisen
må endres.
% vim:fileencoding=utf-8
% -*- encoding: utf-8 -*-
\section*{FORMÅL OG ORGANISASJON}


\subsection*{§ 1}

Realistforeningen er fakultetsforening ved Det
matematisk-naturvitenskapelige fakultet, Universitetet i Oslo.


\subsection*{§ 2}

Realistforeningens formål er å samle foreningens medlemmer til å
drøfte spørsmål av almen og faglig art, å øke kameratskapet mellom
studentene og styrke miljøet ved fakultetet ved selskapelige samvær,
ekskursjoner og andre tiltak.


\subsection*{§ 3}

Realistforeningen har tre kategorier medlemmer: 

\begin{itemize}
\item [a)] Semesterbetalende medlemmer 
\item [b)] Livsvarige medlemmer 
\item [c)] Æresmedlemmer 
\end{itemize}

Semesterbetalende medlemmer og livsvarige medlemmer kan de bli som
studerer, har studert eller underviser ved Det
matematisk-naturvitenskapelige fakultet. Medlemskontingent for
livsvarige medlemmer er 10 x ordinær semesterkontingent.
Æresmedlemmer utnevnes på generalforsamling med 3/4 flertall. Alle
kategorier medlemmer har samme rettigheter.


\subsection*{§ 4}

Hovedstyret avgjør tvilsspørsmål om hvem som er berettiget til
medlemsskap.  Hovedstyret har anledning til å ekskludere et medlem,
som forser seg mot foreningens lover eller viser utillatelig oppførsel
på foreningens arrangementer.


\subsection*{§ 5}

Går en generalforsamling inn for å oppløse foreningen, skal
desisjonsutvalget tidligst tre uker og senest fire måneder senere
innkalle til ekstraordinær generalforsamling hvor saken skal behandles
på nytt. Dersom foreningen oppløses, disponerer generalforsamlingen
over foreningens aktiva.


\subsection*{§ 6}

Realistforeningen har følgende faste styrer og utvalg: 

\begin{itemize}
\item[a)] Hovedstyret 

\item[b)] Arrangementstyret 

\item[c)] Fagstyret

\item[d)] Blæstgruppa

\item[e)] Kjellerstyret

\item[f)] Økonomiutvalget 

\item[g)] Desisjonsutvalget 

\item[h)] Revisjonsutvalget 

\item[i)] Regi 

\item[j)] Vaktgruppa 

\item[k)] Husbjørnredaksjonen
\end{itemize}

I tillegg kan foreningen ha en eller flere komiteer til å ta seg av
oppgaver som faller utenfor de faste styrer og utvalgs virksomhet.



\section*{STYRER OG UTVALG}


\subsection*{§ 7}

Generalforsamlingen vedtar budsjettrammer alle styrer, utvalg og
komiteer skal følge.  Detaljerte budsjettforslag skal utarbeides av
Realistforeningens styrer, utvalg og komiteer og leveres Hovedstyret
innen en måned etter generalforsamlingen.  Alle budsjetter skal være i
tråd med budsjettrammene og godkjennes av Hovedstyret.  Overskridelser
av budsjettrammer i de respektive styrer, utvalg og komiteer behandles
av Økonomiutvalget.

\subsection*{§ 8}

Alle faste styrer og utvalg nevnt i §6, tillitsvalgte valgt av Generalforsamlingen og innehavere av verv
oppnevnt av Hovedstyret ihht. §9e), har ansvar for å lage og oppdatere erfaringsprotokoller
for opplæring av etterfølgende
innehavere av vervet og funksjonærer tilknyttet
styret ihht. §9f). Alle personer omfattet av første ledd har ansvar for at påtroppende vervinnehaver får den
nødvendige opplæring for å inneha sitt verv.
I denne sammenheng skal også disse tekster samt RFs lover gjennomgås.

\subsection*{§ 9 HOVEDSTYRET}

\begin{itemize}

\item[a)] Hovedstyret har følgende stemmeberettigede medlemmer: Formann,
Sekretær, Arrangementssjef, BlæStsjef, Fagsjef, Kjellermester, Regiformann
og Økonomiutvalgets leder.

\item[b)] Hovedstyrets stemmeberettigede velges på generalforsamling for ett
semester. Unntaket er Regiformann som velges ihht. § 17 a).  Alle medlemmene velges ved særskilt valg og må ha vært med
i minst et av Realistforeningens styrer og utvalg nevnt i § 6 minst
ett semester tidligere eller ha innehatt verv nevnt i § 9 e).

\item[c)] Sekretæren er Formannens stedfortreder og fører referat fra
Hovedstyrets møter.

\item[d)] Hovedstyrets oppgave er å koordinere langsiktig virksomhet,
godkjenne program, vedta budsjetter, utarbeide forslag til budsjettrammer 
for Generalforsamlingen og representere Realistforeningen utad.

\item[e)] Hovedstyret oppnevner alle verv. Som verv er å regne:
Skjenkemester, Utlånsansvarlig, Vaktmester, Popvitsjef, Bedriftskontakt,
Donaldsjef, Vaktgruppesjef, Arrangementsmester, Internansvarlig,
IT-ansvarlig og redaktør for Husbjørnen. Hovedstyret kan oppnevne
andre verv etter behov.

\item[f)] Hovedstyret kan bestemme at andre styrer, utvalg og komiteer selv kan
knytte til seg inntil et bestemt antall funksjonærer for å utføre
nærmere bestemte oppgaver.  Funksjonærer blir regnet som medlemmer
av gjeldende styrer og utvalg, dog uten stemmerett.

\item[g)] Innkalling til Hovedstyremøte med foreløpig dagsorden skal være
skriftlig og offentliggjort minst tre dager før møtet, med mindre det
er spesielle tungtveiende hensyn. Innkallingen skal offentliggjøres på
et lett tilgjengelig sted som Hovedstyret bestemmer. Alle RFs
medlemmer kan foreslå saker til dagsorden frem til møtet starter.
Hovedstyret vedtar endelig dagsorden etter godkjenning av innkalling,
og etter det kan bare et enstemmig Hovedstyre endre dagsorden.

\item[h)] Hovedstyret er vedtaksdyktig når der er korrekt innkalt til møtet
og minst 5 av representantene med stemmerett er til stede og saken
vedtas med alminnelig flertall. Ved stemmelikhet teller formannens
stemme dobbelt. Det skal føres møteprotokoll fra hvert av møtene som
godkjennes ved påfølgende hovedstyremøte eller etter at alle
tilstedeværende representanter med stemmerett i Hovedstyret har
mottatt og godkjent protokollen.

\item[i)] Hovedstyret skal innen en måned etter generalforsamlingen vedta en
instruks som spesifiserer arbeidsoppgavene og ansvarsområdene for
foreningens styrer, utvalg og verv.

\end{itemize}


\subsection*{§ 10 ARRANGEMENTSSTYRET}

\begin{itemize}
\item[a)] Arrangementsstyret ledes av Arrangementssjef, og har følgende andre
medlemmer: Regiformann, BlæStsjef, Vaktgruppesjef,
Arrangementsmestere og Kjellermester.

\item[b)] Arrangementsstyret er ansvarlig for planlegging, koordinering og 
gjennomføring av arrangement av sosial karakter. 

\item[c)] Arrangementsstyret skal fremlegge forslag til foreningens sosiale
program for Hovedstyret.

\item[d)] Arrangementsstyret har ansvaret for å rapportere regnskapsrelevant informasjon til Økonomiutvalget og overholde vedtatte budsjetter.
\end{itemize}


\subsection*{§ 11 FAGSTYRET}

\begin{itemize}
\item[a)] Fagstyret ledes av Fagsjef, og har følgende andre medlemmer:
Populærvitenskapelig sjef og Bedriftskontakt.

\item[b)] Fagstyret har ansvar for foreningens faglige profil og kontakt med
næringslivet.

\item[c)] Fagstyret skal fremlegge forslag til foreningens faglige program
for Hovedstyret.

\item[d)] Fagstyret har ansvaret for å rapportere regnskapsrelevant informasjon til Økonomiutvalget og overholde vedtatte budsjetter.
\end{itemize}


\subsection*{§ 12 BLÆSTGRUPPA}

\begin{itemize}
\item[a)] Blæstgruppa ledes av Blæstsjef, og kan knytte til seg så
  mange medlemmer som Blæstsjef finner ønskelig
 
\item[b)] Blæstgruppa har ansvar for å markedsføre alle foreningens
interne og eksterne arrangementer, drive rekruttering og vedlikeholde
foreningens IT-utstyr og nettsider.

\item[c)] Blæstgruppa har ansvaret for å rapportere regnskapsrelevant informasjon til Økonomiutvalget og overholde vedtatte budsjetter.
\end{itemize}


\subsection*{§ 13 KJELLERSTYRET}

\begin{itemize}
\item[a)] Kjellerstyret ledes av Kjellermester, og har følgende andre
medlemmer: Kafesjef, Utlånsansvarlig, Skjenkemestere, Vaktmester.

\item[b)] Kjellerstyret har ansvaret for den daglige drift av
RF-kjellern. Dette omfatter også utlån og vedlikehold av lokalene.

\item[c)] Kjellerstyret har ansvar for å rapportere regnskapsrelevant informasjon til Økonomiutvalget og overholde vedtatte budsjetter.
\end{itemize}


\subsection*{§ 14 ØKONOMIUTVALGET}

\begin{itemize}
\item[a)] Økonomiutvalget har fem medlemmer. Det velges to
  ordinære medlemmer på den
ordinære generalforsamling i hvert semester, og funksjonstiden er to
semestre.  Økonomiutvalgets leder blir hvert semester valgt av
Generalforsamlingen blant utvalgets fire ordinære medlemmer. I
  tillegg velges Forretningsfører for Regi i vårsemesteret, og
  sitter i to semstre. 

\item[b)] Økonomiutvalgets medlemmer kan ikke samtidig være medlemmer av noen
andre av de faste styrer, utvalg eller komiteer nevnt i §6, eller
medlemmer av Bjørnegildestyret. Unntaket er Økonomiutvalgets leder,
som er medlem i Hovedstyret, og Forretningsfører for Regi, som er
medlem i Regi.

\item[c)] Økonomiutvalget har ansvaret for Realistforeningens regnskap og
for å lære opp alle styrer, utvalg og komiteer i økonomistyring.
Økonomiutvalget kan pålegge styrer, utvalg og komiteer å føre sine
egne regnskap, men fører ellers alle regnskap.  Økonomiutvalget skal
også kontrollere at foreningens budsjett blir fulgt.  I tilfelle
budsjettsprekk, skal det aktuelle styret, det aktuelle utvalget
eller den aktuelle komiteen samt Hovedstyret, Revisjonsutvalget og
Desisjonsutvalget informeres.

\item[d)] Økonomiutvalget skal avholde konstituerende møte innen 10 virkedager
etter nyvalg. Her velger utvalget en sekretær som skal føre protokoll
over alle møter.  Det konstituerende møtet skal innkalles av lederen i
det fungerende Økonomiutvalget og samtlige medlemmer av dette
innkalles.  Det sittende Økonomiutvalg fører regnskapene ut den 
inneværende periode.

\item[e)] Revisjonsutvalget og ett medlem av Hovedstyret, i tillegg til
Økonomiutvalgets leder, har møte-, tale-, og forslagsrett på
Økonomiutvalgets møter. Utvalget kan pålegge medlemmer av styrer og
komiteer å møte ved behandlingen av bestemte saker.

\item[f)] Økonomiutvalget kan bare fatte vedtak i møte når det er minst tre
medlemmer tilstede. For gyldig vedtak kreves det at minst tre
medlemmer har stemt for forslaget. Ved stemmelikhet teller leders
stemme dobbelt.

\item[g)] Økonomiutvalget skal utarbeide forskrifter som kan lette kontrollen
med regnskapene.

\item[h)] Økonomiutvalget overtar driften av Realistforeningen inntil nytt
Hovedstyre er valgt dersom det sittende ikke kan funksjonere. Ingen
utbetalinger, med unntak av utestående fordringer, skal skje før en
generalforsamling er avholdt.
\end{itemize}
 

\subsection*{§ 15 DESISJONSUTVALGET}

\begin{itemize}
\item[a)] Desisjonsutvalget har tre medlemmer: Det velges ett medlem på den
ordinære generalforsamlingen i hvert semester, og funksjonstiden er
tre semestre.

\item[b)] Valgbare er alle som har hatt valgte verv som nevnt i §21 l).

\item[c)] Desisjonsutvalgets medlemmer kan ikke samtidig være medlemmer av
noen andre av de faste styrer, utvalg eller komiteer nevnt i §6, eller
verv nevnt i §9e eller medlemmer av Bjørnegildestyret.

\item[d)] Desisjonsutvalget har den endelige avgjørelse i tvilsspørsmål om
tolkning av lovene. Utvalget kan også fatte vedtak og gi regler i
situasjoner hvor lovene måtte vise seg å være utilstrekkelige. Ethvert
medlem av Realistforeningen har rett til å innanke for
Desisjonsutvalget vedtak hvor det kan være tvil om lovligheten.

\item[e)] Ved mistanke om misligheter kan Desisjonsutvalget suspendere
medlemmer av styrer, utvalg og komiteer. Suspensjonen kan omfatte et
organ i sin helhet, selv om det ikke foreligger konkret mistanke mot
hvert enkelt medlem.  I tilfelle suspensjon er foretatt skal
Desisjonsutvalget straks sørge for at det blir innkalt til
ekstraordinær generalforsamling der mistillitsforslag behandles og
nyvalg eventuelt avholdes.

\item[f)] Desisjonsutvalget har ansvar for at Realistforeningens arkiv til
enhver tid er i orden.
        
\item[g)] Desisjonsutvalget har møte- og talerett i alle foreningens organer. 

\item[h)] Desisjonsutvalget har ansvar for at lovtekstene oppdateres og er tilgjengelige.
\end{itemize}


\subsection*{§ 16 REVISJONSUTVALGET}

\begin{itemize}
\item[a)] Minst ett medlem til Revisjonsutvalget velges på generalforsamling
hvert semester, slik at Revisjonsutvalget til enhver tid har tre
medlemmer. Funksjonstiden er to semestre.

\item[b)] Revisjonsutvalgets medlemmer kan ikke samtidig være medlemmer av
Hovedstyret, Økonomiutvalget, Desisjonsutvalget eller Bjørnegildestyret 
eller inneha vervet som Forretningsfører i Regi eller verv som
medfører stemmerett i Kjellerstyret eller Arrangementstyret, eller ha
vært medlem av Økonomiutvalget foregående to semestre. To medlemmer av et 
styre, en komité eller et utvalg nevnt i §6 kan ikke samtidig delta i 
revideringen av et regnskap.

\item[c)] Revisjonsutvalgets oppgave er å revidere Realistforeningens
regnskaper. Minst to av revisjonsutvalgets medlemmer må delta i
revideringen av et regnskap.

\item[d)] Alle regnskaper skal være innlevert senest tre uker før
generalforsamling påfølgende semester. Blir ikke regnskapene godkjent
på generalforsamlingen, skal Økonomiutvalget inndra alle bilag og
fullføre regnskapet. Det kan gis dispensasjon til avvik fra dette
punkt av Hovedstyret i samarbeid med Revisjons- og Økonomiutvalget.

\item[e)] På Generalforsamlingen skal Revisjonsutvalget legge fram
revisjonsberetningen, som kan være skrevet av Revisjonsutvalget selv
eller, hvis Hovedstyret finner det nødvendig, en av Kredittilsynet
godkjent revisor.  Revisjonsutvalget har ansvar for å opplyse
Generalforsamlingen om eventuelle budsjettoverskridelser.
\end{itemize}


\subsection*{§ 17 REGI}

\begin{itemize}
\item[a)] Regiformann og Forretningsfører velges på generalforsamling for ett
år; Regiformann om høsten, Forretningsfører om våren.  Øvrige
medlemmer godkjennes av Hovedstyret etter innstilling fra Regiformann.

\item[b)] Regis oppgave er å stå for drift og forvaltning av
Realistforeningens tekniske utstyr.

\item[c)] Regi har ansvaret for å rapportere regnskapsrelevant
  informasjon til Økonomiutvalget og overholde vedtatte budsjetter.
\end{itemize}


\subsection*{§ 18 VAKTGRUPPA}

\begin{itemize}
\item[a)] Vaktgruppa ledes av Vaktgruppesjef som utpekes av Hovedstyret etter
innstilling av Vaktgruppa.

\item[b)] Vaktgruppas oppgave er i samråd med Arrangementstyret og
Kjellerstyret å stå for vakthold under Realistforeningens
arrangementer.

\item[c)] Vaktgruppa har ansvar for Vaktgruppas regnskap og økonomistyring.
\end{itemize}


\subsection*{§ 19 HUSBJØRNREDAKSJONEN}
\begin{itemize}
\item[a)] Husbjørnredaksjonen ledes av Husbjørnredaktør, og kan knytte
  til seg så mange medlemmer som Husbjørnredaktør finner ønskelig.

\item[b)] Husbjørnredaksjonen har ansvar for å utgi avisen
  Husbjørnen, \emph{Ursus Domesticus}.

\item[c)] Husbjørnredaksjonen har ansvar for å rapportere
  regnskapsrelevant informasjon til Økonomiutvalget og overholde
  vedtatte budsjetter. 
\end{itemize}

\subsection*{§ 20 BJØRNEGILDET} 

\begin{itemize}
\item[a)] Bjørnegildet avholdes vårsemesteret hvert tredje år. 

\item[b)] Bjørnegildet ledes av et styre, hvis medlemmer velges av
generalforsamlingen seneste tre semestere før Bjørnegildet.

\item[c)] Bjørnegildestyret består av Gildesjef, Sekretær, Økonomiansvarlig
og så mange medlemmer som generalforsamlingen finner nødvendig.

\item[d)] Bjørnegildestyret har ansvaret for Bjørnegildets regnskap og
økonomistyring.

\item[e)] Formannen i Realistforeningen har møte- og stemmerett i
Bjørnegildestyret.
\end{itemize}


\subsection*{§ 21 KOMITEER}

\begin{itemize}
\item[a)] Komiteer utnevnes og får sitt mandat av en generalforsamling eller
Hovedstyret.

\item[b)] En komité har ansvar for å rapportere regnskapsrelevant informasjon til Økonomiutvalget og overholde vedtatte budsjetter.
\end{itemize}


\subsection*{§ 22 GENERALFORSAMLING}

\begin{itemize}
\item[a)] Generalforsamlingen er foreningens høyeste myndighet i spørsmål som
ikke kommer inn under §15 d) under første punkt. Generalforsamlingen er
beslutningsdyktig når minst 1/10 av medlemmene er tilstede, dog slik
at 50 stemmeberettigede er tilstrekkelig dersom foreningen har mer enn
500 medlemmer.

\item[b)] Ordinær generalforsamling avholdes i andre halvdel av hvert 
semester.  Ekstraordinær generalforsamling avholdes når Hovedstyret 
vedtar det eller det kreves av Desisjonsutvalget eller minst 1/10 av 
medlemmene, dog slik at 50 medlemmer er tilstrekkelig dersom foreningen 
har mer enn 500 medlemmer.

\item[c)] Innkallelse til ordinær og ekstraordinær generalforsamling må være
offentliggjort minst 10 virkedager i forveien.  Ved ordinær og ekstraordinær
generalforsamling må forslag til foreløpig dagsorden være offentliggjort 
senest 5 virkedager i forveien.  Som virkedag regnes alle dager i samme 
semeseter som ikke er helg, offentlig høytidsdag eller feriedag for 
studentene ved Det matematisk- naturvitenskapelige fakultet ved 
Universitetet i Oslo.  Generalforsamlinger innkalles av Hovedstyret.  
Dersom dette ikke fungerer eller ikke etterkommer lovlige krav om at 
generalforsamling skal kalles inn, skal Desisjonsutvalget overta 
Hovedstyrets plikter når det gjelder generalforsamlinger, med unntak av 
§22 e).

\item[d)] Forslag om lovendringer og andre saker som ønskes tatt opp på
generalforsamlingen må være Hovedstyret i hende og offentliggjøres 5
virkedager før.  Desisjonsutvalget kan fremme endringsforslag inntil 48
timer før generalforsamlingen.  Lovendringsforslag kan ikke behandles
på ekstraordinær generalforsamling.

\item[e)] Hovedstyret skal offentliggjøre et sett med budsjettrammer til
Generalforsamlingens overveielse senest 10 virkedager før høstens
Generalforsamling.  Andre forslag til budsjettrammer som ønskes
vedtatt må være Hovedstyret i hende og offentliggjøres senest 48
timer før Generalforsamlingen.

\item[f)] Generalforsamlingen kan foreta endringer i rekkefølgen av punktene
til det endelige forslag til dagsorden.  Den kan også utelukke ett
eller flere av de foreslåtte punkter så lenge det ikke strider mot
§22 l).  Den endelige dagsorden godkjennes av generalforsamlingen.  I
forbindelse med godkjennelse av dagsorden skal det velges ordstyrer,
referent og to medlemmer til å undertegne generalforsamlingens
protokoll.

\item[g)] Ethvert medlem kan på generalforsamlingen foreslå tatt opp saker
utenom den oppsatte dagsorden. Generalforsamlingen kan ikke fatte
vedtak i slike saker.

\item[h)] Generalforsamlingen kan med alminnelig flertall gi ikke-medlemmer
møte- og talerett.

\item[i)] Avstemningen på generalforsamlinger skal være hemmelig når minst
tre stemmeberettigede krever det.  Stemmerett har alle medlemmer av
Realistforeningen som var innmeldt før innkallingen til
generalforsamlingen.

\item[j)] Valgbare til verv i Realistforeningen er alle foreningens
medlemmer, med de innskrenkninger som følger av §14 b), §15 b) og c) og 
§16 b).

\item[k)] Valg på flere tillitsvalgte under ett avgjøres med alminnelig
flertall.  Ved valg på en enkelt tillitsvalgt kan tre
stemmeberettigede kreve at valget skal avgjøres med absolutt
flertall.  Oppnår ingen dette ved første avstemming, avholdes bundet
omvalg.

\item[l)] På ordinær generalforsamling behandles:

        \begin{itemize}
        \item[1.] Regnskaper, etter en redegjørelse for RFs totale økonomi. 
        
        \item[2.] Budsjettrammer.  På høstens generalforsamling vedtas 
                  budsjettrammer for neste år.  På vårens generalforsamling 
                  kan budsjettrammene revideres.

        \item[3.] Eventuelle lovendringsforslag 

        \item[4.] Semesterberetninger 

        \item[5.] Fastsettelse av kontingenter. 

        \item[6.] Valg av tillitsvalgte: 

                \begin{itemize}
                \item[6a)] Formann (§9b) 
    
                \item[6b)] Sekretær (§9b) 

                \item[6c)] Arrangementsjef (§9b) 

                \item[6d)] Fagsjef (§9b) 

                \item[6e)] Kjellermester(§9b) 

                \item[6f)] BlæStsjef (§9b) 

                \item[6g)] Kafesjef (§13a) 

                \item[6h)] To Skjenkemestre (§13a)

                \item[6i)] To medlemmer til Økonomiutvalget (§14a)

                \item[6j)] Leder av Økonomiutvalget (§14a) 

                \item[6k)] Ett medlem til Desisjonsutvalget (§15a) 

                \item[6l)] Minst ett medlem til Revisjonsutvalget (§16a) 

                \item[6m)] Formann i Regi (§17a) 

                \item[6n)] Forretningsforfører i Regi (§17a) 
                \end{itemize}
        \end{itemize}
\end{itemize}


\subsection*{§ 23 MISTILLIT}

Foreningens medlemmer kan fremme mistillitsforslag mot medlemmer av
styrer, utvalg og komiteer unntatt Desisjonsutvalget. Slike forslag
kan bare behandles av en generalforsamling, og må være fremmet 48
timer før generalforsamlingen. Mistillitsforslag vedtas med 2/3
flertall.  Dersom mistillitsforslaget mot et medlem av et organ blir
vedtatt, kan generalforsamlingen vedta å holde nyvalg på samtlige
medlemmer av organet for resten av hvert medlems funksjonstid.


\subsection*{§ 24 LOVENDRINGER}

Forslag til lovendring skal bare behandles på ordinær
generalforsamling, og må få 2/3 flertall blant de tilstedeværende
stemmeberettigede for å vedtas.


\subsection*{§ 25 LOVERS GYLDIGHET}

Disse lovene er gyldige fra den dag de blir vedtatt, slik at alle
tidligere lover opphører å gjelde fra samme dag.





\chapter{Takk til}
Alle som har vært med å skrive denne boka, alle som kommer til å oppdatere den i fremtiden, og alle andre som har gjort, gjør eller kommer til å gjøre en innsats for RF.

\chapter{Endringslogg}
\input{changelog}
\end{document}
